\SourcePage{188}
\SGAsection{3}{同调代数补充}
本节的主要目的是定义前面宣布的\Term{isomorphism}{同构}
\[
\operatorname{Biext}^{1}(P,Q;G)\simeq\operatorname{Ext}^{1}(P\mathbin{\otimes^{\mathrm L}}Q,G).
\]
为此，我们需要作一些本身也有独立意义的\Term{homological-algebra}{同调代数}方面的讨论，这将占据 3.1 至 3.4 各小节。3.1 和 3.3 的内容可视为 [4] 的补充，尤其是该文 6.9 节的补充。

\SGAsection{3.1}{由阿贝尔范畴中的链复形构造加性余纤维范畴}
设 $\mathcal C$ 为一个\Term{abelian-category}{阿贝尔范畴}，$L_\bullet$ 为 $\mathcal C$ 中的一个\Term{chain-complex}{链复形}，
\[
L_2\longrightarrow L_1\longrightarrow L_0\longrightarrow0,\tag{3.1.1}
\]
我们采用下标\Term{degree}{次数}的约定来写它（\Term{differential}{微分算子}的\Term{degree}{次数}为 $-1$）。因而 $i<0$ 时 $L_i=0$。我们将定义 $\mathcal C$ 上的一个\Term{cofibered-category}{余纤维范畴} $\Psi=\Psi_{L_\bullet}$（在\Term{cofibered-category}{余纤维范畴}的\Term{canonical-equivalence}{典范等价}意义下，它就是上述文献 6.9 中的同名\Term{category}{范畴}）。$\Psi_{L_\bullet}$ 的\Term{object}{对象}是二元组
\[
(X,\alpha),
\]
\SourcePage{189}
其中 $X$ 是 $L_0$ 以 $\mathcal C$ 的某个\Term{object}{对象} $A$（依赖于 $X$）为核的\Term{extension}{扩张}：
\[
0\longrightarrow A\xrightarrow{i}X\xrightarrow{j}L_0\longrightarrow0,\tag{3.1.2}
\]
而 $\alpha$ 是此\Term{extension}{扩张}沿 $d_0:L_1\to L_0$ 所得\Term{inverse-image}{逆像} $d_0^*(X)$ 的一个\Term{trivialization}{平凡化}，并要求相应的、沿 $d_1:L_2\to L_1$ 所得\Term{inverse-image}{逆像} $d_1^*(d_0^*(X))$ 的\Term{trivialization}{平凡化}，是由\Term{transitivity-isomorphism}{传递同构}
\[
d_1^*(d_0^*(E))\simeq(d_0d_1)^*(E)
\]
及关系 $d_0d_1=0$ 给出的显然\Term{trivialization}{平凡化}。还可以换一种方式描述满足上述条件的数据 $\alpha$：它是 $d_0:L_1\to L_0$ 到 $\alpha:L_1\to X$ 的一个\Term{lifting}{提升}，满足 $\alpha d_1=0$：
\[
\begin{tikzcd}[column sep=large,row sep=small]
&&0\arrow[d]\\
&&A\arrow[d,"i"]\\
&&X\arrow[d,"j"]\\
L_2\arrow[r,"d_1"']&L_1\arrow[r,"d_0"']\arrow[ur,"\alpha"]&L_0\arrow[d]\\
&&0
\end{tikzcd}
\qquad\alpha d_1=0.\tag{3.1.3}
\]
在这种\Term{form}{形式}下可以看出，\Term{object}{对象} $(X,\alpha)$ 的描述只依赖于 $L_\bullet$ 的一阶\Term{truncation}{截断}\Term{chain-complex}{链复形} $[L_\bullet]$：它由 $L_\bullet$ 通过“消去 $i\geq2$ 时的 $H_i(L_\bullet)$”得到，其\Term{degree}{次数} $\geq2$ 的\Term{set-component}{分量}为零，\Term{degree}{次数} $\leq0$ 的\Term{set-component}{分量}与 $L_\bullet$ 相同，而\Term{degree}{次数} $1$ 的\Term{set-component}{分量}是 $d_1:L_2\to L_1$ 的\Term{cokernel}{余核}。从 $(X,\alpha)$ 到 $(X',\alpha')$ 的\Term{homomorphism}{同态}定义为与 $L_0$ 上的恒等相容、且满足 $u\alpha=\alpha'$ 的\Term{extension}{扩张}\Term{homomorphism}{同态} $u:X\to X'$；\footnote{DELIGNE 指出，描述 $\Psi_{L_\bullet}(A)$ 的一种更简单方式，是把它看作 $L_\bullet$ 以 $A$（视为只在零次非零的\Term{complex}{复形}）为核的\Term{complex}{复形}\Term{extension}{扩张}所成的\Term{category}{范畴}。}
\SourcePage{190}
这些\Term{homomorphism}{同态}的复合由相应的\Term{extension}{扩张}\Term{homomorphism}{同态}的复合给出。这样就得到了所说的\Term{category}{范畴} $\Psi_{L_\bullet}$。在\Term{canonical-isomorphism}{典范同构}意义下，它只依赖于 $L_\bullet$ 的一阶\Term{truncation}{截断} $[L_\bullet]$。

有一个显然的\Term{functor}{函子} $(X,\alpha)\mapsto A$：
\[
\pi:\Psi_{L_\bullet}\longrightarrow\mathcal C.\tag{3.1.4}
\]
这个\Term{functor}{函子}是\Term{cofibered-functor}{余纤维化函子}。事实上，设 $u:A\to B$ 是 $\mathcal C$ 中的一个\Term{morphism}{态射}，$(X,\alpha)$ 是 $A$ 处的\Term{fiber-category}{纤维范畴} $\Psi_{L_\bullet}(A)$ 的一个\Term{object}{对象}。那么 $X$ 是 $L_0$ 以 $A$ 为核的\Term{extension}{扩张}；可以考虑它沿 $u$ 的\Term{direct-image}{顺像}，即 $L_0$ 以 $B$ 为核的\Term{extension}{扩张} $u_*(X)$，并赋予它复合的部分\Term{trivialization}{平凡化} $\beta=f\alpha$，其中 $f:X\to u_*(X)$ 是\Term{canonical}{典范}\Term{morphism}{态射}。立即可验证
\[
f:(X,\alpha)\longrightarrow(u_*(X),\beta)
\]
是\Term{functor}{函子} (3.1.4) 的一个\Term{cocartesian-morphism}{余笛卡儿态射}，即它对 $\Psi_{L_\bullet}$ 中以 $(X,\alpha)$ 为源的 $u$-\Term{morphism}{态射}满足通常的\Term{universal-property}{泛性质}；为此，只需使用\Term{morphism}{态射} $f:X\to u_*(X)$ 的类似\Term{universal-property}{泛性质}，后者是众所周知且显然的。

最后，\Term{cofibered-category}{余纤维范畴} $\Psi_{L_\bullet}$ 是\Term{additive}{加性的}，即对 $\mathcal C$ 的任意两个\Term{object}{对象} $A,B$，它满足下面两个条件：
\[
\Psi_{L_\bullet}(0)\text{ 与单点范畴等价},\tag{3.1.5}
\]
\[
\Psi_{L_\bullet}(A\times B)\longrightarrow
\Psi_{L_\bullet}(A)\times\Psi_{L_\bullet}(B)\text{ 是等价}.\tag{3.1.6}
\]
利用 $L_0$ 以 $\mathcal C$ 中变动\Term{object}{对象}为核的\Term{extension}{扩张}所成\Term{cofibered-category}{余纤维范畴}的相应性质，验证是显然的。我们将自由使用
\SourcePage{191}
$\Psi_{L_\bullet}$ 的\Term{additivity}{加性}的各种推论；它们已在上述文献第 1 节中建立，并在前面（2.5）回顾过。

\textbf{3.1.7。} 这些推论之一是：对 $\mathcal C$ 的每个\Term{object}{对象} $A$，$\Psi_{L_\bullet}(A)$ 中有一个\Term{zero-object}{零对象}，或称\Term{trivial-object}{平凡对象}（2.5.1）。容易验证，它由 $L_0$ 以 $A$ 为核的\Term{trivial-extension}{平凡扩张} $X=L_0+A$ 给出，赋予它的部分\Term{trivialization}{平凡化}由\Term{set-component}{分量}为 $d_0$ 和 $0$ 的 $\alpha:L_1\to X$ 定义。另一方面（2.5.2），对 $\Psi_{L_\bullet}(A)$ 的每个\Term{object}{对象} $(X,\alpha)$，其\Term{automorphism-group}{自同构群}都\Term{canonically-isomorphic}{典范同构}于 $\Psi_{L_\bullet}(A)$ 的\Term{trivial-object}{平凡对象}的\Term{automorphism-group}{自同构群}；将它记为 $\Psi^0_{L_\bullet}(A)$，这是一个\Term{commutative-group}{交换群}。当然，从 $\Psi_{L_\bullet}(A)$ 的具体描述中也可直接得到这一事实，甚至得到一个\Term{canonical-isomorphism}{典范同构}
\[
\begin{aligned}
\Psi^0_{L_\bullet}(A)
&=\operatorname{Ker}\bigl(\operatorname{Hom}(L_0,A)\longrightarrow\operatorname{Hom}(L_1,A)\bigr)\\
&\simeq\operatorname{Hom}(H_0(L_\bullet),A)\simeq\operatorname{Ext}^0(L_\bullet,A),
\end{aligned}\tag{3.1.8}
\]
它由 $X$ 的\Term{extension}{扩张}\Term{automorphism-group}{自同构群}与\Term{group}{群} $\operatorname{Hom}(L_0,A)$ 之间通常的\Term{isomorphism}{同构}诱导。\Term{isomorphism}{同构} (3.1.7.1) 显然关于 $A$ 具有\Term{functoriality}{函子性}。

\textbf{3.1.9。} 最后回顾（2.5.2），借助\Term{fiber-category}{纤维范畴} $\Psi_{L_\bullet}(A)$ 中的自然\Term{multiplication}{乘法}结构，它的\Term{isomorphism-class}{同构类}\Term{set}{集合}自然带有\Term{group}{群}结构；我们将其记为 $\Psi^1_{L_\bullet}(A)$。后者关于 $A$ 是一个\Term{additive-functor}{加性函子}。

\SourcePage{192}
\SGAsection{3.2}{确定 $\Psi^1_{L_\bullet}(A)$}
考虑 $\Psi_{L_\bullet}(A)$ 的一个\Term{object}{对象} $(X,\alpha)$（3.1.3）。考虑由\Term{morphism}{态射} $j:X\to L_0$ 定义的、\Term{length}{长度}为 $1$ 的\Term{cochain-complex}{上链复形}
\[
[X\longrightarrow L_0]:\quad0\longrightarrow X\longrightarrow L_0\longrightarrow0\cdots,\tag{3.2.1}
\]
其零次\Term{set-component}{分量}为 $X$。这是 $A$ 的一个\Term{resolution}{消解}，因而作为\Term{derived-category}{导出范畴} $D(\mathcal C)$ [5] 的\Term{object}{对象}，它可以与 $A$ 等同。于是图 (3.1.3) 给出一个\Term{degree}{次数}为 $1$ 的\Term{complex}{复形}\Term{homomorphism}{同态}
\[
L_\bullet\longrightarrow[X\longrightarrow L_0],\tag{3.2.2}
\]
具体由下面各方块\Term{anticommutative}{反交换}的图给出：
\[
\begin{tikzcd}[column sep=large]
\cdots\arrow[r]&L_2\arrow[r,"d_1"]&L_1\arrow[r,"d_0"]\arrow[d,"u^{-1}=-\alpha"']&L_0\arrow[r]\arrow[d,"u^0=\mathrm{id}"]&0\arrow[r]&\cdots\\
\cdots\arrow[r]&0\arrow[r]&X\arrow[r,"j"']&L_0\arrow[r]&0\arrow[r]&\cdots
\end{tikzcd}\tag{3.2.3}
\]
（这里采用通常的约定：把\Term{chain-complex}{链复形} $L_\bullet$ 对应于\Term{differential}{微分算子}\Term{degree}{次数}为 $1$、第 $i$ 次\Term{set-component}{分量}为 $L_{-i}$ 的\Term{complex}{复形}。）这样便在\Term{derived-category}{导出范畴}中得到一个\Term{degree}{次数}为 $1$ 的\Term{homomorphism}{同态}
\[
c(X,\alpha):L_\bullet\longrightarrow A
\qquad\text{（$D(\mathcal C)$ 中次数为 $1$ 的同态）},\tag{3.2.3}
\]
它显然只依赖于 $(X,\alpha)$ 的\Term{isomorphism-class}{同构类}。由此定义了\Term{canonical-map}{典范映射}
\[
\Psi^1_{L_\bullet}(A)\longrightarrow\operatorname{Ext}^1(L_\bullet,A)
\overset{\mathrm{def}}=\operatorname{Hom}_{D(\mathcal C)}(L_\bullet,A[1]).\tag{3.2.4}
\]
\SourcePage{193}
\begin{SGAresult}{\Term{theorem}{定理}}{3.2.5}
映射 (3.2.4) 是\Term{group}{群}\Term{isomorphism}{同构}，并关于 $\mathcal C$ 的\Term{object}{对象} $A$ 具有\Term{functoriality}{函子性}。
\end{SGAresult}
\begin{proof}
我们将映射 (3.2.4) 的\Term{additivity}{加性}及其关于 $A$ 的\Term{functoriality}{函子性}留给读者，只证明它是\Term{bijection}{双射}。对于\Term{injectivity}{单射性}，先证明：若 (3.2.3) 的\Term{complex}{复形}\Term{homomorphism}{同态} $u$ 在 $D(\mathcal C)$ 中为零，则它在模\Term{homotopy}{同伦}的\Term{complex}{复形}\Term{category}{范畴} $K(\mathcal C)$ 中已经为零；换言之，存在\Term{morphism}{态射} $s:L_0\to X$ 使得
\[
js=u^0=\mathrm{id}_{L_0},\qquad sd_0=-u^{-1}=\alpha.\tag{3.2.5.1}
\]
事实上，假设的含义是：将所考虑的\Term{complex}{复形}\Term{morphism}{态射}与适当的\Term{quasi-isomorphism}{拟同构} $v$ 复合后，结论便成立：
\[
\begin{tikzcd}
0\arrow[r]&X\arrow[r,"j"]\arrow[d,"v^0"']&L_0\arrow[r]\arrow[d,"v^1"]&0\\
0\arrow[r]&C^0\arrow[r,"j'"']&C^1\arrow[r]&C^2\arrow[r]&\cdots.
\end{tikzcd}\tag{*}
\]
此外，可以假定 $i\geq2$ 时 $C^i=0$（只需将 $i\geq2$ 的 $C^i$ 换成 $0$，并将 $C^1$ 换成 $\operatorname{Im}(C^0\to C^1)=\operatorname{Ker}(C^1\to C^2)$）。但此时，\Term{morphism}{态射} (*) 在 $A$ 上诱导恒等，故它把 $X$ 等同于 $L_0$ 和 $C^0$ 在 $C^1$ 上的\Term{fiber-product}{纤维积}；而给定一个定义\Term{homotopy}{同伦} $vu\sim0$、因而满足 $j's'=v^1u^0$ 的\Term{morphism}{态射} $s':L_0\to C^0$，便得到一个\Term{morphism}{态射} $s:L_0\to X$，使 $js=u^0=\mathrm{id}_{L_0}$。这样，$s$ 定义了 $L_0$ 以 $A$ 为核的\Term{extension}{扩张} $X$ 的一个\Term{splitting}{分裂}，故\Term{complex}{复形} $X\to L_0$ 在 $K(\mathcal C)$ 中\Term{isomorphic}{同构}于只在 $A$ 处非零的\Term{complex}{复形} $K^\bullet$。
\SourcePage{194}
对于后者，显然\Term{homomorphism}{同态}
\[
\operatorname{Hom}_{K(\mathcal C)}(L_\bullet,K^\bullet)
\longrightarrow\operatorname{Hom}_{D(\mathcal C)}(L_\bullet,K^\bullet)
\]
是\Term{isomorphism}{同构}，由此即得我们的断言。

既然如此，给定一个满足 (3.2.5.1) 的\Term{homomorphism}{同态} $j:L_0\to X$，恰好\Term{equivalence}{等价}于给定 $L_0$ 以 $A$ 为核的\Term{extension}{扩张} $X$ 的一个与部分\Term{trivialization}{平凡化} $\alpha$ 相容的\Term{trivialization}{平凡化}。这就得到了 (3.2.4) 的\Term{injectivity}{单射性}。

对于\Term{surjectivity}{满射性}，从右端的一个元素出发；按定义，它可以用一个\Term{degree}{次数}为 $1$ 的\Term{complex}{复形}\Term{homomorphism}{同态}描述：
\[
\begin{tikzcd}[column sep=large]
\cdots\arrow[r]&L_2\arrow[r,"d_1"]&L_1\arrow[r,"d_0"]\arrow[d,"u^{-1}"']&L_0\arrow[r]\arrow[d,"u^0"]&0\\
&0\arrow[r]&C^0\arrow[r]&C^1\arrow[r]&\cdots
\end{tikzcd}\tag{3.2.5.2}
\]
此图各方块\Term{anticommutative}{反交换}，$C^\bullet$ 是 $A$ 的一个\Term{resolution}{消解}，并且如前可假定 $i\geq2$ 时 $C^i=0$。于是 $C^0$ 是 $C^1$ 以 $A$ 为核的一个\Term{extension}{扩张}，可考虑其沿 $L_0\to C^1$ 的\Term{inverse-image}{逆像}\Term{extension}{扩张} $X$。图 (3.2.5.2) 中方块的\Term{anticommutativity}{反交换性}表明，$-u^{-1}$ 通过一个\Term{morphism}{态射} $\alpha:L_1\to X$ \Term{resolution}{分解}，它是一个部分\Term{splitting}{分裂}，即满足 $j\alpha=d_0$（$j$ 见 (3.1.2)）。此外，(3.2.5.2) 中的关系 $u^{-1}d_1=0$ 蕴含 $\alpha d_1=0$。故 $(X,\alpha)$ 确为 $\Psi_{L_\bullet}(A)$ 的一个元素。而且，由构造本身，\Term{homomorphism}{同态} (3.2.5.2) 是对应于 $(X,\alpha)$ 的\Term{homomorphism}{同态} (3.2.3) 与下面\Term{homomorphism}{同态}的复合：
\[
[X\longrightarrow L_0]\longrightarrow C^\bullet=[C^0\longrightarrow C^1],
\]
\SourcePage{195}
其\Term{set-component}{分量}为 $(f,u^0)$，其中 $f:X\to C^0$ 是由 $X$ 作为\Term{extension}{扩张} $C^0$ 的\Term{inverse-image}{逆像}的定义所得的\Term{canonical}{典范}\Term{homomorphism}{同态}。上面的\Term{homomorphism}{同态}是 $A$ 的\Term{resolution}{消解}之间的\Term{homomorphism}{同态}（在 $A$ 上诱导恒等），因此由 $\operatorname{Ext}^1=\operatorname{Hom}^1_{D(\mathcal C)}$ 的定义可知，(3.2.5.2) 所描述的给定元素 $\operatorname{Ext}^1(L_\bullet,A)$ 正是 $c(X,\alpha)$。这就证明了 (3.2.4) 的\Term{surjectivity}{满射性}，并完成了 3.2.5 的证明。
\end{proof}

\SGAsection{3.3}{余纤维范畴 $\Psi_{L_\bullet}$ 的性质}
本小节并非理解后文所必需。这里我们自由使用 [4] 中的概念，目的是把上述结果与该文的结果作比较。

\begin{SGAresult}{\Term{theorem}{定理}}{3.3.1}
$\mathcal C$ 上的\Term{additive}{加性}\Term{cofibered-category}{余纤维范畴} $\Psi_{L_\bullet}$ 是\Term{left-exact}{左正合的}（上述文献 2.2），它具有一个\Term{quasi-maximal-object}{拟极大对象}（quasi-maximal；上述文献 4.4），并且对该\Term{category}{范畴}的每个\Term{object}{对象} $E=(X,\alpha)$，都存在 $E$ 的一个\Term{typical-complex}{典型复形}（complexe typique）$L^E_\bullet$（上述文献 3.1）。此外，$\Psi_{L_\bullet}$ 的\Term{typical-pro-complex}{典型 pro-复形}（pro-complexe typique；上述文献 3.7 和 4.4）在 $D(\mathcal C)$ 中\Term{canonically-isomorphic}{典范同构}于 3.1 中考虑的\Term{truncation}{截断}\Term{complex}{复形} $[L_\bullet]$。
\end{SGAresult}
下面具体说明 3.3.1 中各个断言的含义。

\textbf{3.3.2。} 说\Term{cofibered-category}{余纤维范畴} $\Psi_{L_\bullet}$ 是\Term{left-exact}{左正合的}，是指它是\Term{additive}{加性的}，并且还满足条件：

\textbf{(3.3.2.1)} 对 $\mathcal C$ 中每个\Term{exact-sequence}{正合列}
\[
0\longrightarrow A'\xrightarrow{u}A\xrightarrow{v}A''\longrightarrow0,
\]
\SourcePage{196}
自然\Term{functor}{函子}（记号见上述文献 1.8）
\[
\Psi_{L_\bullet}(A)\longrightarrow\Psi_{L_\bullet}([A\longrightarrow A''])
\]
是\Term{equivalence-of-categories}{范畴等价}。

回顾一下：若 $\Psi$ 是 $\mathcal C$ 上的一个\Term{cofibered-category}{余纤维范畴}，$v:A\to A''$ 是 $\mathcal C$ 中的\Term{morphism}{态射}，则 $\Psi([A\to A''])$ 表示由二元组 $(E,\rho)$ 构成的\Term{category}{范畴}，其中 $E$ 是 $\Psi(A)$ 的\Term{object}{对象}，而 $\rho$ 是 $\Psi(A'')$ 中相应\Term{object}{对象} $v_*(E)$ 的一个“\Term{trivialization}{平凡化}”，即该\Term{object}{对象}与 $\Psi(A'')$ 的\Term{zero-object}{零对象}（2.5.1）之间的一个\Term{isomorphism}{同构}。(3.3.2.1) 中\Term{functor}{函子}
\[
\Psi(A')\longrightarrow\Psi([A\longrightarrow A''])
\]
的定义是显然的。在 $\Psi=\Psi_{L_\bullet}$ 的情形，$\Psi_{L_\bullet}([A\to A''])$ 的\Term{object}{对象}具体为三元组 $(X,\alpha,\rho)$，其中 $(X,\alpha)$ 如 (3.1.3)，而 $\rho:X\to A''$ 是 $\mathcal C$ 中的一个\Term{morphism}{态射}，满足
\[
\rho i=v,\qquad\rho\alpha=0.\tag{*}
\]
当 $v:A\to A''$ \Term{embedding}{嵌入}如 3.2.3 中的一个\Term{short-exact-sequence}{短正合列}时，众所周知，给定满足 (*) 中第一个关系的\Term{morphism}{态射}，\Term{equivalence}{等价}于给定 $L_0$ 以 $A'\simeq\operatorname{Ker}v$ 为核的一个\Term{extension}{扩张} $X'$，使它成为 $X$ 的子\Term{extension}{扩张}（取 $X'=\operatorname{Ker}\rho$）；这正是 $L_0$ 以 $\mathcal C$ 中变动\Term{object}{对象}为核的\Term{extension}{扩张}所成\Term{cofibered-category}{余纤维范畴}的\Term{left-exact}{左正合性}。而 (*) 的第二个关系表示 $\alpha$ \Term{resolution}{分解}为一个 $\alpha':L_1\to X'$。除去一些显然可核实的相容性，这就证明了 (3.2.3) 中具体说明的\Term{left-exact}{左正合性}。

\SourcePage{197}
\textbf{3.3.3。} 回顾一下，$\mathcal C$ 上的\Term{cofibered-category}{余纤维范畴} $\Psi$ 的一个\Term{object}{对象} $E_0$ 称为\Term{quasi-maximal-object}{拟极大对象}（quasi-maximal），如果对 $\Psi$ 的每个位于 $\mathcal C$ 的\Term{object}{对象} $A$ 之上的\Term{object}{对象} $E$，均可在 $\Psi$ 中找到一个\Term{morphism}{态射} $E\to E'$，它在 $\mathcal C$ 中诱导一个\Term{monomorphism}{单态射} $A\to A'$，并使 $E$ 被 $E_0$ 所\Term{domination}{支配}，即存在一个\Term{morphism}{态射} $E_0\to E'$。若 $\mathcal C$ 有\Term{enough-injectives}{足够多的内射对象}，这\Term{equivalence}{等价}于说，$\Psi$ 中每个位于 $\mathcal C$ 的一个\Term{injective-object}{内射对象}之上的\Term{object}{对象}，都被 $E_0$ 所\Term{domination}{支配}。在 $\Psi=\Psi_{L_\bullet}$ 的情形，有下面这个\Term{canonical}{典范}的\Term{quasi-maximal-object}{拟极大对象} $E_0$。置
\[
L'_1=\operatorname{Coker}(d_1:L_2\longrightarrow L_1),
\]
以 $E_0$ 表示 $\Psi_{L_\bullet}(L'_1)$ 的\Term{object}{对象} $(X_0,\alpha_0)$，其中 $X_0=L_0+L'_1$ 是 $L_0$ 以 $L'_1$ 为核的\Term{trivial-extension}{平凡扩张}，而 $\alpha_0$ 的\Term{set-component}{分量}为 $(d'_1,\mathrm{id}_{L'_1})$；这里 $d'_1:L_1\to L_0$ 是由 $d_1:L_1\to L_0$ 取商得到的。

为证明 $E_0$ 确为\Term{quasi-maximal-object}{拟极大对象}，注意对 $\Psi_{L_\bullet}$ 的每个\Term{object}{对象} $(X,\alpha)$，都可找到一个从它到某个 $(X',\alpha')$ 的\Term{morphism}{态射}，该\Term{morphism}{态射}在 $\mathcal C$ 中诱导一个\Term{monomorphism}{单态射} $A\to A'$，并使 $L_0$ 以 $A'$ 为核的\Term{extension}{扩张} $X'$ \Term{split-extension}{分裂}：事实上，只需取 $A'=X$、$(X',\alpha')=i_*(X,\alpha)$，其中记号如 (3.1.2)。另一方面，若 $\Psi$ 的一个\Term{object}{对象} $(X',\alpha')$ 所对应的\Term{extension}{扩张} $X'$ \Term{split-extension}{分裂}，则它显然被 $E_0$ 所\Term{domination}{支配}。

\textbf{3.3.4。} 设 $E$ 为 $\mathcal C$ 上的一个\Term{additive}{加性}\Term{cofibered-category}{余纤维范畴} $\Psi$ 的\Term{object}{对象}，位于 $\mathcal C$ 的\Term{object}{对象} $A$ 之上。考虑关于 $\mathcal C$ 中变动\Term{object}{对象} $B$ 的\Term{functor}{函子}
\[
B\longmapsto\operatorname{Hom}(E,\theta_B),\tag{3.3.4.1}
\]
\SourcePage{198}
其中 $\theta_B$ 表示 $\Psi(B)$ 的\Term{trivial-object}{平凡对象}。若上述\Term{functor}{函子}可由 $\mathcal C$ 的某个\Term{object}{对象} $\Omega_E$ \Term{representable-functor}{表示}，则称 $E$ 具有一个\Term{typical-complex}{典型复形}（complexe typique）。所考虑的\Term{functor}{函子}\Term{canonical}{典范}地映入由 $A$ \Term{representable-functor}{表示}的\Term{functor}{函子} $B\mapsto\operatorname{Hom}(A,B)$，因此得到一个\Term{canonical}{典范}\Term{homomorphism}{同态}
\[
A\longrightarrow\Omega_E,\tag{3.3.4.2}
\]
而 $\mathcal C$ 中相应的\Term{length}{长度}为 $1$ 的\Term{chain-complex}{链复形} $A\to\Omega_E$，称为 $E$ 的\Term{typical-complex}{典型复形}。

在 $\Psi=\Psi_{L_\bullet}$、$E=(X,\alpha)$ 的情形，立即可验证，\Term{functor}{函子} (3.3.4.1) 由 $\operatorname{Coker}\alpha$ \Term{representable-functor}{表示}，而\Term{canonical}{典范}\Term{homomorphism}{同态} (3.3.4.2) 就是复合
\[
A\xrightarrow{i}X\xrightarrow{\mathrm{can}}\operatorname{Coker}\alpha.
\]

\textbf{3.3.5。} 若 $\mathcal C$ 上的一个\Term{additive}{加性}\Term{cofibered-category}{余纤维范畴} $\Psi$ 具有一个\Term{quasi-maximal-object}{拟极大对象}（3.3.3），且其各\Term{object}{对象}的\Term{typical-complex}{典型复形}均存在，则可证明：$\Psi$ 的一个\Term{quasi-maximal-object}{拟极大对象} $E_0$ 的\Term{typical-complex}{典型复形} $T_\bullet$，作为\Term{derived-category}{导出范畴} $D(\mathcal C)$ 的\Term{object}{对象}，在\Term{canonical-isomorphism}{典范同构}意义下不依赖于 $E_0$ 的选择（上述文献 4.4）。它称为\Term{additive}{加性}\Term{cofibered-category}{余纤维范畴} $\Psi$ 的\Term{typical-complex}{典型复形}。最后，若进一步假定 $\Psi$ 是\Term{left-exact}{左正合的}（3.3.2），则可证明：由\Term{typical-complex}{典型复形} $T_\bullet$ 可以重构 $\Psi$，精确到 $\mathcal C$ 上\Term{fibered-category}{纤维范畴}的\Term{equivalence-of-categories}{等价}（该\Term{equivalence-of-categories}{等价}在唯一\Term{isomorphism}{同构}意义下唯一；上述文献 6.10），并有关于 $\mathcal C$ 的\Term{object}{对象} $A$ 具有\Term{functoriality}{函子性}的\Term{canonical-isomorphism}{典范同构}（上述文献 5.4）：
\[
\Psi^0(A)\xrightarrow{\sim}\operatorname{Ext}^0(T_\bullet,A),\qquad
\Psi^1(A)\xrightarrow{\sim}\operatorname{Ext}^1(T_\bullet,A).\tag{3.3.5.1}
\]
\SourcePage{199}
为完成与上述文献的联系，我们来证明：在 $\Psi=\Psi_{L_\bullet}$ 的情形，$\Psi_{L_\bullet}$ 的\Term{typical-complex}{典型复形}是\Term{truncation}{截断} $[L_\bullet]$。这样便通过 (3.3.5.1) 重新给出\Term{isomorphism}{同构} (3.1.8) 和 (3.2.4) 的定义；可以验证，它与 3.1 和 3.2 中的定义相容。

为证明此断言，显然可以只考虑 $L_\bullet$ 已经一阶\Term{truncation}{截断}的情形（因为在 3.1 中已指出，\Term{truncation}{截断}并不改变\Term{category}{范畴} $\Psi_{L_\bullet}$，精确到\Term{cofibered-category}{余纤维范畴}的\Term{isomorphism}{同构}）。此时，由 3.3.4 给出的\Term{typical-complex}{典型复形}的具体构造，显然 3.3.3 中描述的\Term{quasi-maximal-object}{拟极大对象}的\Term{typical-complex}{典型复形}\Term{canonically-isomorphic}{典范同构}于 $L_\bullet$。

\textbf{注记 3.3.6。} 设
\[
0\longrightarrow A'\longrightarrow A\longrightarrow A''\longrightarrow0
\]
为 $\mathcal C$ 中的一个\Term{exact-sequence}{正合列}，由此得到关于 $\operatorname{Ext}^i$ 的\Term{exact-sequence}{正合列}：
\[
0\longrightarrow\operatorname{Ext}^0(L_\bullet,A')\longrightarrow\operatorname{Ext}^0(L_\bullet,A)
\longrightarrow\cdots\longrightarrow\operatorname{Ext}^1(L_\bullet,A)\longrightarrow\operatorname{Ext}^1(L_\bullet,A'').
\]
（引入 $i\geq2$ 的 $\operatorname{Ext}^i$ 后，还可向右延续。）利用\Term{isomorphism}{同构} (3.1.8) 和 (3.2.4)，便得到\Term{exact-sequence}{正合列}
\[
\begin{aligned}
0&\longrightarrow\Psi^0_{L_\bullet}(A')\longrightarrow\Psi^0_{L_\bullet}(A)
\longrightarrow\Psi^0_{L_\bullet}(A'')\xrightarrow{\partial}\Psi^1_{L_\bullet}(A')\\
&\longrightarrow\Psi^1_{L_\bullet}(A)\longrightarrow\Psi^1_{L_\bullet}(A'').
\end{aligned}\tag{3.3.6.1}
\]
其中除标作 $\partial$ 的\Term{arrow}{箭头}外，其余\Term{arrow}{箭头}都由 $\Psi^0_{L_\bullet}$ 和 $\Psi^1_{L_\bullet}$ 的\Term{functoriality}{函子性}直接定义。另一方面，上述文献 (2.5.2) 对 $\mathcal C$ 上任意\Term{left-exact}{左正合}\Term{cofibered-category}{余纤维范畴} $\Psi$，借助一个\Term{homomorphism}{同态}
\SourcePage{200}
\[
\partial:\Psi^0(A'')\longrightarrow\Psi^1(A')\tag{3.3.6.2}
\]
定义了\Term{exact-sequence}{正合列} (3.3.6.1)：它把左端的元素 $u$ 对应到右端中由 $\Psi(A')$ 的某个元素给出的类；此元素在\Term{canonical-isomorphism}{典范同构}意义下由以下条件确定：它在 $\Psi([A\to A''])$（记号见 3.3.2）中的像是二元组 $(X_0,u)$，其中 $X_0$ 是 $\Psi(A)$ 的\Term{trivial-object}{平凡对象}。现在我断言，在我们所关心的 $\Psi=\Psi_{L_\bullet}$ 的情形，用这种“几何”方式定义的\Term{homomorphism}{同态} (3.3.6.2)，正是 (3.3.6.1) 中的\Term{boundary-homomorphism}{边界同态} $\partial$ 改变\Term{signature}{符号}所得。因此，作出上述等同并改变 $\partial$ 的\Term{signature}{符号}后，后一个\Term{exact-sequence}{正合列}就是上述文献 (2.5.2) 中的\Term{exact-sequence}{正合列}。考虑到 3.3.5 中指出的相容性，这实际上包含在上述文献 5.4 的最后一个断言中。

\textbf{注记 3.3.6.3。} 借此机会，我们来澄清 [4，第 3.11 节] 中定义基本\Term{isomorphism}{同构} (3.11.3) 和 (3.11.4) 时一个不清楚之处（这两个\Term{isomorphism}{同构}在该文 5.4 的证明中起着关键\Term{action}{作用}）。该文第 35 页的图是交换的，而非\Term{anticommutative}{反交换}的，因此它本身并不定义一个\Term{homomorphism}{同态}
\[
u:[L^X_1\longrightarrow L^X_0]\longrightarrow K^\bullet[1].
\]
应将其中的 $u^1$ 换成 $-u^1$，使图\Term{anticommutative}{反交换}。（也可以采用相反约定，将 $u^0$ 换成 $-u^0$；不过，上述约定，即把 $u^1$ 改成 $-u^1$，才使所得\Term{isomorphism}{同构} (3.11.3) 如该文所宣布的那样，推广该文的\Term{isomorphism}{同构} (3.5.3)。）这样澄清定义后，还须修正该文 5.4 的断言：那里所考虑的两个\Term{coboundary-homomorphism}{上边界同态}
\SourcePage{201}
相差一个\Term{signature}{符号}。

\textbf{3.3.7。} 在上述文献 6.13 中，还说明了当 $L_\bullet$ 变动时，$\mathcal C$ 上\Term{cofibered-category}{余纤维范畴} $\Psi_{L_\bullet}$ 的\Term{variance}{变性}；尤其由\Term{chain-complex}{链复形} $L_\bullet$ 和 $L'_\bullet$ 重构了从 $\Psi_{L_\bullet}$ 到 $\Psi_{L'_\bullet}$ 的\Term{cocartesian-morphism}{余笛卡儿态射}所成的\Term{category}{范畴}。特别地，从 $\Psi_{L_\bullet}$ 到 $\Psi_{L'_\bullet}$ 的\Term{cocartesian-functor}{余笛卡儿函子}的\Term{isomorphism-class}{同构类}\Term{set}{集合}，与\Term{isomorphism}{同构}的\Term{set}{集合} $\operatorname{Hom}_{D(\mathcal C)}(L'_\bullet,L_\bullet)$ 之间有一个\Term{canonical}{典范}的一一对应。我们只使用这些信息中的一小部分，并在下一小节中利用 3.1 所给出的 $\Psi_{L_\bullet}$ 的简单描述加以展开。

\SGAsection{3.4}{$\Psi_{L_\bullet}$ 关于 $L_\bullet$ 的变性}
考虑一个\Term{complex}{复形}\Term{morphism}{态射}
\[
f:L'_\bullet\longrightarrow L_\bullet,\tag{3.4.1}
\]
它由\Term{commutative-diagram}{交换图}
\[
\begin{tikzcd}
\cdots\arrow[r]&L'_2\arrow[r]\arrow[d,"f_2"']&L'_1\arrow[r]\arrow[d,"f_1"']&L'_0\arrow[r]\arrow[d,"f_0"]&0\\
\cdots\arrow[r]&L_2\arrow[r]&L_1\arrow[r]&L_0\arrow[r]&0
\end{tikzcd}\tag{3.4.2}
\]
给出。我们要给它对应一个\Term{canonical}{典范}\Term{functor}{函子}
\[
f^*:\Psi_{L_\bullet}\longrightarrow\Psi_{L'_\bullet},\tag{3.4.3}
\]
它实际上是 $\mathcal C$ 上\Term{cofibered-category}{余纤维范畴}之间的一个\Term{cocartesian-functor}{余笛卡儿函子}（即与取值于 $\mathcal C$ 的结构\Term{functor}{函子}相容，并把\Term{cocartesian-morphism}{余笛卡儿态射}变成\Term{cocartesian-morphism}{余笛卡儿态射}；参见 SGA 1 第 VI 讲）。若 $(X,\alpha)$ 是 $\Psi_{L_\bullet}$ 的一个\Term{object}{对象}，则以 $f^*(X,\alpha)$ \Term{representing-object}{表示对象}
\SourcePage{202}
$(X',\alpha')$，其中 $X'$ 是给定的、$L_0$ 以 $A$ 为核的\Term{extension}{扩张} $X$ 沿 $f_0$ 的\Term{inverse-image}{逆像}，因而是 $L'_0$ 以 $A$ 为核的\Term{extension}{扩张}；$\alpha'$ 则是借助 (3.4.2) 中第二个方块的\Term{transitivity}{传递性}，从部分\Term{trivialization}{平凡化} $\alpha$ 得到的部分\Term{trivialization}{平凡化}。（$\alpha'$ 与关系 $d'_0d'_1=0$ 的相容条件，显然由 $\alpha$ 的类似条件得到。）由构造可直接看出，\Term{functor}{函子} (3.4.3) 与取值于 $\mathcal C$ 的结构\Term{functor}{函子}交换。它是\Term{cocartesian-functor}{余笛卡儿的}，则立即由下述事实得出：在 $\mathcal C$ 的\Term{object}{对象}的\Term{extension}{扩张}\Term{category}{范畴}中，“\Term{extension}{扩张}的\Term{direct-image}{顺像}”运算与“\Term{extension}{扩张}的\Term{inverse-image}{逆像}”运算交换。

注意，\Term{functor}{函子} (3.4.3) 定义了\Term{functor}{函子}之间的\Term{homomorphism}{同态}
\[
f^i:\Psi^i_{L_\bullet}\longrightarrow\Psi^i_{L'_\bullet},\qquad i=0,1\tag{3.4.4}
\]
（比较 2.6）。另一方面，\Term{homomorphism}{同态} (3.4.1) 定义了\Term{homomorphism}{同态}
\[
f^i:\operatorname{Ext}^i(L_\bullet,-)\longrightarrow\operatorname{Ext}^i(L'_\bullet,-),\qquad i\in\mathbb Z.\tag{3.4.5}
\]
现在我断言，$i=0,1$ 时，\Term{homomorphism}{同态} (3.4.4) 和 (3.4.5) 与形如 (3.1.8) 和 (3.2.4) 的\Term{isomorphism}{同构}相容，即对 $A\in\operatorname{ob}\mathcal C$、$i=0,1$，下图交换：
\[
\begin{tikzcd}
\Psi^i_{L_\bullet}(A)\arrow[r,"\sim"]\arrow[d]&\operatorname{Ext}^i(L_\bullet,A)\arrow[d]\\
\Psi^i_{L'_\bullet}(A)\arrow[r,"\sim"]&\operatorname{Ext}^i(L'_\bullet,A).
\end{tikzcd}\tag{3.4.6}
\]
\SourcePage{203}
这种相容性的验证基本上是显然的，留给读者。

\begin{SGAresult}{命题}{3.4.7}
\Term{functor}{函子} (3.4.3) 是\Term{faithful}{忠实的}（相应地，\Term{fully-faithful}{充分忠实的}，或为\Term{equivalence-of-categories}{范畴等价}）的充要条件是：$H_0(f):H_0(L'_\bullet)\to H_0(L_\bullet)$ 是\Term{epimorphism}{满态射}（相应地，$H_0(f)$ 是\Term{isomorphism}{同构}且 $H_1(f)$ 是\Term{epimorphism}{满态射}，或 $H_0(f)$ 和 $H_1(f)$ 都是\Term{isomorphism}{同构}）。
\end{SGAresult}
\begin{proof}
事实上，以 $f^i$ 表示 (3.4.4) 中具有\Term{functoriality}{函子性}的\Term{homomorphism}{同态}（$i=0,1$），所需条件是 $f^0$ 为\Term{monomorphism}{单态射}（相应地，$f^0$ 是\Term{isomorphism}{同构}且 $f^1$ 为\Term{monomorphism}{单态射}，或 $f^0,f^1$ 都是\Term{isomorphism}{同构}）。由 (3.4.6) 的\Term{commutativity}{交换性}，这些条件可用 $\operatorname{Ext}^i$、$i=0,1$ 这两个\Term{functor}{函子}上的类似条件来表达。此时断言只是\Term{homological-algebra}{同调代数}中一个相当显然的命题，我们把它作为练习留给读者。
\end{proof}

\textbf{3.4.8。} 现在考虑一个\Term{chain-complex}{链复形}的\Term{exact-sequence}{正合列}
\[
0\longrightarrow L_\bullet\xrightarrow{f}L'_\bullet\xrightarrow{g}L''_\bullet\longrightarrow0,\tag{3.4.8.1}
\]
由此得到 (3.4.3) 型的\Term{functor}{函子}
\[
\Psi_{L''_\bullet}\xrightarrow{g^*}\Psi_{L'_\bullet}\xrightarrow{f^*}\Psi_{L_\bullet}.
\]
对 (3.4.3) 型\Term{functor}{函子}可以预见的\Term{transitivity}{传递性}，在这里给出一个\Term{functor}{函子}\Term{isomorphism}{同构}
\[
f^*g^*\simeq(gf)^*=(0_{L_\bullet\to L''_\bullet})^*.
\]
\SourcePage{204}
最后一项显然\Term{canonically-isomorphic}{典范同构}于从 $\Psi_{L''_\bullet}$ 到 $\Psi_{L_\bullet}$ 的“零”\Term{cocartesian-functor}{余笛卡儿函子}，因此上述\Term{isomorphism}{同构}可用于定义一个\Term{canonical}{典范}\Term{functor}{函子}
\[
\Psi_{L''_\bullet}\longrightarrow\Psi_{[L_\bullet\longrightarrow L'_\bullet]},\tag{3.4.8.2}
\]
其中 $\Psi_{[L_\bullet\to L'_\bullet]}$ 表示由二元组 $(X',t)$ 构成的\Term{category}{范畴}，$X'$ 是 $\Psi_{L'_\bullet}$ 的\Term{object}{对象}，$t$ 是其\Term{inverse-image}{逆像} $f^*(X')$ 的一个“\Term{trivialization}{平凡化}”，即一个\Term{isomorphism}{同构} $f^*(X')\xrightarrow{\sim}0_{L_\bullet,A}$；这里 $0_{L_\bullet,A}$ 表示 $\Psi_{L_\bullet}(A)$ 的\Term{zero-object}{零对象}，$A$ 是 $X'$ 之下的 $\mathcal C$ 的\Term{object}{对象}。现在我断言，\Term{functor}{函子} (3.4.8.2) 是\Term{equivalence-of-categories}{范畴等价}（这以表达式 $\Psi_{L_\bullet}(A)$ 关于 $L_\bullet$ 的一种\Term{exactness}{正合性}，补充了 3.3.1 所建立的关于 $A$ 的\Term{left-exact}{左正合性}）。

下面勾勒这一断言的证明。给定 (3.4.8.2) 右端的一个\Term{object}{对象}，\Term{equivalence}{等价}于给定 $L'_0$ 以 $\mathcal C$ 中某个\Term{object}{对象} $A$ 为核的一个\Term{extension}{扩张}
\[
0\longrightarrow A\xrightarrow{i'}X'\xrightarrow{j'}L'_0\longrightarrow0,
\]
再加上一个满足 $j'\alpha'=d'_0$、$\alpha'd'_1=0$ 的\Term{morphism}{态射} $\alpha':L'_1\to X'$，以及一个满足 $j't=f_0$、$td_0=\alpha'f_1$ 的\Term{morphism}{态射} $t:L_0\to X'$。参见下图（虚线部分为下面构造的\Term{object}{对象}和\Term{morphism}{态射}）：
\[
\begin{tikzcd}[column sep=large,row sep=large]
&&0\arrow[d]&0\arrow[d,dashed]\\
&&A\arrow[d,"i'"]&A\arrow[d,dashed,"i''"]\\
&&X'\arrow[d,"j'"]\arrow[r,dashed]&X''\arrow[d,dashed,"j''"]\\
0\arrow[r]&L_0\arrow[r,"f_0"]\arrow[ur,bend left=20,"t"]&L'_0\arrow[r,"g_0"]&L''_0\arrow[r]&0\\
0\arrow[r]&L_1\arrow[u,"d_0"]\arrow[r,"f_1"]&L'_1\arrow[u,"d'_0"]\arrow[uu,bend left=30,"\alpha'"]\arrow[r,"g_1"]&L''_1\arrow[u,"d''_0"]\arrow[uu,bend right=30,dashed,"\alpha''"']\arrow[r]&0\\
0\arrow[r]&L_2\arrow[u,"d_1"]\arrow[r]&L'_2\arrow[u,"d'_1"]\arrow[r]&L''_2\arrow[u,"d''_1"]\arrow[r]&0
\end{tikzcd}
\qquad
\left\{\begin{aligned}
j'\alpha'&=d'_0,\\
\alpha'd'_1&=0,\\
j't&=f_0,\\
td_0&=\alpha'f_1.
\end{aligned}\right.
\]
\SourcePage{205}
将一个\Term{object}{对象} $L$ 以另一个\Term{object}{对象} $A$ 为核的\Term{extension}{扩张}关于第一个变元的众所周知的\Term{left-exact}{左正合性}应用于图的第一行，可知给定满足 $j't=f_0$ 的 $t$，就能把\Term{extension}{扩张} $X'$ 解释为 $L''_0$ 以 $A$ 为核的某个\Term{extension}{扩张} $X''$ 的\Term{inverse-image}{逆像}，并使 $t$ 对应于该\Term{extension}{扩张}沿 $g_0f_0=0$ 所得\Term{inverse-image}{逆像}的自然\Term{trivialization}{平凡化}。给定满足 $j'\alpha'=d'_0$ 和 $\alpha'f_1=td_0$ 的 $\alpha'$，也可通过基本相同的对应，只是应用于图的第二行，解释为来自某个满足 $j''\alpha''=d''_0$ 的\Term{morphism}{态射} $\alpha'':L''_1\to X''$（将 $\alpha'$ 和 $\alpha''$ 解释为\Term{extension}{扩张} $X'$ 和 $X''$ 分别沿 $d'_0$ 和 $d''_0$ 所得\Term{inverse-image}{逆像}的\Term{trivialization}{平凡化}）。最后，关系 $\alpha'd'_1=0$ 此时\Term{equivalence}{等价}于 $\alpha''d''_1=0$。总之，给定 $\Psi_{[L_\bullet\to L'_\bullet]}$ 的\Term{object}{对象} $(X',\alpha',t)$，\Term{equivalence}{等价}于给定 $\Psi_{L''_\bullet}$ 的\Term{object}{对象} $(X'',\alpha'')$；这就基本证明了 (3.4.8.2) 确为\Term{equivalence-of-categories}{范畴等价}的断言。

\textbf{3.4.9。} 继续考虑\Term{exact-sequence}{正合列} (3.4.8.1)，以及它定义的 $\operatorname{Ext}^i$ 的\Term{exact-sequence}{正合列}：
\[
\begin{aligned}
0&\longrightarrow\operatorname{Ext}^0(L''_\bullet,A)\xrightarrow{g^0}\operatorname{Ext}^0(L'_\bullet,A)
\xrightarrow{f^0}\operatorname{Ext}^0(L_\bullet,A)\\
&\xrightarrow{\partial^0}\operatorname{Ext}^1(L''_\bullet,A)\xrightarrow{g^1}\operatorname{Ext}^1(L'_\bullet,A)
\xrightarrow{f^1}\operatorname{Ext}^1(L_\bullet,A)\longrightarrow\cdots.
\end{aligned}\tag{3.4.9.1}
\]
(3.1.8) 和 (3.2.4) 型的\Term{canonical-isomorphism}{典范同构}使我们能用\Term{category}{范畴} $\Psi_{L_\bullet}$、$\Psi_{L'_\bullet}$ 和 $\Psi_{L''_\bullet}$，对这个\Term{exact-sequence}{正合列}的前六项作“几何”解释；相容性 (3.4.6) 则使我们同样能对该图中的\Term{arrow}{箭头}
\SourcePage{206}
$f^0,g^0,f^1,g^1$ 作几何解释。我们也将利用刚才证明的结果，即 (3.4.8.2) 为\Term{equivalence-of-categories}{范畴等价}，对该图的\Term{arrow}{箭头} $\partial^0$ 作几何解释。事实上，设
\[
u\in\operatorname{Ext}^0(L_\bullet,A),\qquad\text{即 }u:L_0\longrightarrow A,\quad ud_0=0.
\]
考虑 $\Psi_{[L_\bullet\to L'_\bullet]}$ 的\Term{object}{对象} $(0,\bar u)$，其中 $0$ 是 $\Psi_{L'_\bullet}$ 的\Term{zero-object}{零对象}，故其沿 $f_0$ 的\Term{inverse-image}{逆像}与 $\Psi_{L_\bullet}$ 的\Term{zero-object}{零对象}等同；而 $\bar u$ 是借助\Term{canonical-isomorphism}{典范同构} (3.1.8) 由 $u$ 定义的后一个\Term{object}{对象}的\Term{automorphism}{自同构}。因为 (3.4.8.2) 是\Term{equivalence-of-categories}{等价}，上述\Term{object}{对象} $(0,\bar u)$ 来自 $\Psi_{L''_\bullet}(A)$ 的一个在唯一\Term{isomorphism}{同构}意义下确定的\Term{object}{对象}。这个\Term{object}{对象}的\Term{isomorphism-class}{同构类}是 $\Psi^1_{L''_\bullet}(A)\simeq\operatorname{Ext}^1(L''_\bullet,A)$ 的一个元素。我断言，这个元素正是 $\partial^0(u)$。

我们将把这种相容性归结于 [4，推论 5.5] 中的类似相容性（其中的\Term{signature}{符号}应按 3.3.6.3 所述修正），从而证明它；这一论证也将给出\Term{equivalence-of-categories}{等价} (3.4.8.2) 的另一个证明，即把它归结于类似的陈述 [4，命题 2.7]。为此，考虑 $\mathcal C$ 的\Term{opposite-category}{反范畴} $\mathcal C^\circ$，并用\Term{category}{范畴} $\mathcal C^\circ$ 以及 $\mathcal C^\circ$ 上的\Term{cofibered-category}{余纤维范畴} $\mathcal E$ 来解释固定 $A$ 时 $\Psi_{L_\bullet}(A)$ 的构造，以及 3.4.8 和本小节 3.4.9 中的构造；这里 $\mathcal E$ 是 $\mathcal C^\circ$ 的固定\Term{object}{对象} $A$ 以 $\mathcal C^\circ$ 的变动\Term{object}{对象} $L$ 为核的\Term{extension}{扩张}所成的\Term{cofibered-category}{余纤维范畴}。这是 $\mathcal C^\circ$ 上的一个\Term{additive}{加性}\Term{cofibered-category}{余纤维范畴}，甚至在 $\mathcal C^\circ$ 上\Term{left-exact}{左正合}（按上述文献 2.2 的意义）。$\mathcal C$ 中的一个\Term{chain-complex}{链复形} $L_\bullet$ 可以解释为
\SourcePage{207}
$\mathcal C^\circ$ 中的一个\Term{cochain-complex}{上链复形}；在此解释下，所给出的 $\Psi_{L_\bullet}(A)$ 的构造，与上述文献 1.8 意义下 $\widehat{\mathcal E}(L_\bullet)$ 的\Term{opposite-category}{反范畴}的构造相同。\Term{functor}{函子} (3.4.8.2) 此时等同于上述文献 (2.1.2) 型\Term{functor}{函子}的反\Term{functor}{函子}，其中 $\mathcal E$ 被替换为它在\Term{cochain-complex}{上链复形}\Term{category}{范畴}上的延拓 $\widehat{\mathcal E}$；故该\Term{functor}{函子}为\Term{equivalence-of-categories}{等价}这一事实，就是上述文献 2.7 a)$\Rightarrow$d) 的特例。另一方面，\Term{exact-sequence}{正合列} (3.4.9.1) 与 $\operatorname{Ext}^i(A,L''_\bullet)$、$\operatorname{Ext}^i(A,L'_\bullet)$ 和 $\operatorname{Ext}^i(A,L_\bullet)$ 的\Term{exact-sequence}{正合列}等同，只是 $\partial^i$ 相差\Term{signature}{符号} $(-1)^i$（Verdier 说的）。同样，\Term{isomorphism}{同构} (3.1.8) 和 (3.2.4) 与类似的\Term{isomorphism}{同构} (3.11.5) 和 (3.11.4) 等同。因此，所陈述的相容性是上述文献 5.5（按前述方式修正后）的特例。

\textbf{注记 3.4.9.2。} 整理者很怕\Term{signature}{符号}，坦言自己对前面的证明并非完全有把握，因为担心从 $\mathcal C$ 转到 $\mathcal C^\circ$ 时遗漏了\Term{signature}{符号}。他建议有勇气的读者直接验证所宣布的相容性（或者相反的相容性）。请注意，这些\Term{signature}{符号}问题在后面的各讲中（幸好）不会起任何\Term{action}{作用}。

\SGAsection{3.5}{任意 $\mathbb Z$-模的一个典范平坦部分消解}
我们将把上述结果应用于下述情形：$\mathcal C$ 是\Term{topos}{拓扑斯} $\mathcal T$ 的\Term{commutative-group}{交换群}所成的\Term{category}{范畴}，即 $\mathcal T$ 的 $\mathbb Z$-\Term{module}{模}所成的\Term{category}{范畴}。设 $P$ 是 $\mathcal C$ 的一个\Term{object}{对象}，即\Term{topos}{拓扑斯} $\mathcal T$ 的一个\Term{commutative-group}{交换群}。我们将定义一个关于 $P$ 具有\Term{functoriality}{函子性}、向 $P$ \Term{augmentation}{增广}的\Term{chain-complex}{链复形} $L_\bullet(P)$：
\SourcePage{208}
\[
\begin{tikzcd}[column sep=small]
0\arrow[r]&\underset{L_2}{\mathbb Z[P\times P]+\mathbb Z[P\times P\times P]}\arrow[r,"d_1"]&\underset{L_1}{\mathbb Z[P\times P]}\arrow[r,"d_0"]&\underset{L_0}{\mathbb Z[P]}\arrow[r]\arrow[d,"\epsilon"]&0\\
&&&P&
\end{tikzcd}\tag{3.5.1}
\]
满足 $i\geq3$ 时 $L_i(P)=0$。这个\Term{complex}{复形}的\Term{set-component}{分量}，已在上面的 (3.5.1) 中用 1.4 的记号具体给出。为具体说明\Term{differential}{微分算子} $d_0,d_1$ 及\Term{augmentation}{增广} $\epsilon$，我们用 $[p]$ 表示由 $P$ 的点 $p\in P(S)$ 定义的 $\mathbb Z[P](S)$ 的点；类似地，当 $p,q,r\in P(S)$ 时，用 $[p,q]$、$[p,q,r]$ 分别表示 $\mathbb Z[P\times P](S)$ 和 $\mathbb Z[P\times P\times P](S)$ 的元素。采用这些记号，用下式定义 $\epsilon,d_0,d_1$：
\[
\left\{\begin{aligned}
\epsilon[p]&=p,\\
d_0[p,q]&=[p+q]-[p]-[q],\\
d_1[p,q]&=[p,q]-[q,p],\\
d_1[p,q,r]&=[p+q,r]-[p,q+r]+[p,q]-[q,r].
\end{aligned}\right.\tag{3.5.2}
\]
容易验证，这样定义的\Term{commutative-group}{交换群}\Term{homomorphism}{同态}确实定义了一个向 $P$ \Term{augmentation}{增广}的\Term{complex}{复形} $L_\bullet(P)$，即有关系
\[
\epsilon d_0=0,\qquad d_1d_0=0.
\]
事实上，第一个关系是显然的，而第二个关系按 $L_2(P)$ 的两个因子分成两个关系；它们分别表达了 $P$ 上给定\Term{composition-law}{合成律}的\Term{commutativity}{交换性}与\Term{associativity}{结合性}。
\SourcePage{209}
这样定义的\Term{augmented-complex}{增广复形}显然关于 $P$ 具有\Term{functoriality}{函子性}：一个\Term{morphism}{态射} $u:P\to P'$ 给出\Term{commutative-diagram}{交换图}
\[
\begin{tikzcd}
L_\bullet(P)\arrow[r,"L_\bullet(u)"]\arrow[d]&L_\bullet(P')\arrow[d]\\
P\arrow[r,"u"']&P'.
\end{tikzcd}
\]
还要注意，\Term{complex}{复形} $L_\bullet$ 的各\Term{set-component}{分量}都是\Term{flat}{平坦的}，正如任意\Term{free-module}{自由} $\mathbb Z$-\Term{module}{模}一样。

\begin{SGAresult}{命题}{3.5.3}
有 $H_1(L_\bullet(P))=0$，且\Term{augmentation}{增广} $\epsilon$ 诱导一个\Term{isomorphism}{同构} $H_0(L_\bullet(P))\xrightarrow{\sim}P$。
\end{SGAresult}
\begin{proof}
为证明 3.5.3，可用标准方法归结到 $\mathcal T$ 是\Term{set}{集合}\Term{category}{范畴}的情形；这时 $P$ 就是一个普通的 $\mathbb Z$-\Term{module}{模}，而 (3.5.1) 很可能是“众所周知的”。我们将给出一个不使用这一归约的直接证明：把 3.4.7 应用于\Term{augmentation}{增广}\Term{morphism}{态射}
\[
\epsilon:L_\bullet(P)\longrightarrow P,
\]
其中照例将 $P$ 视为仅在零次非零的\Term{chain-complex}{链复形}。于是只需证明，对 $\mathcal C$ 的每个\Term{object}{对象} $G$，\Term{functor}{函子}
\[
\epsilon^*:\Psi_P(G)\longrightarrow\Psi_{L_\bullet(P)}(G)\tag{*}
\]
是\Term{equivalence-of-categories}{范畴等价}。而 (*) 左端显然
\SourcePage{210}
\Term{isomorphic}{同构}于 $\mathbb Z$-\Term{module}{模} $P$ 以 $G$ 为核的\Term{extension}{扩张}所成的\Term{category}{范畴}。为具体说明右端，借助 1.4，将 $L_0,L_1,L_2$ 以 $G$ 为核的\Term{extension}{扩张}所成的\Term{category}{范畴}，分别解释为 $G_P$-\Term{torsor}{挠子}的\Term{category}{范畴}、$G_{P\times P}$-\Term{torsor}{挠子}的\Term{category}{范畴}，以及由一个 $G_{P\times P}$-\Term{torsor}{挠子}和一个 $G_{P\times P\times P}$-\Term{torsor}{挠子}组成的二元组的\Term{category}{范畴}。因此，(*) 右端的一个\Term{object}{对象}可等同于二元组 $(E,\varphi)$，其中 $E$ 是一个 $G_P$-\Term{torsor}{挠子}，而 $\varphi$ 是 $P\times P$ 上一个\Term{torsor}{挠子} $E'$ 的\Term{trivialization}{平凡化}；$E'$ 对应于由 $E$ 定义的 $\mathbb Z(P)$ 以 $G$ 为核的\Term{extension}{扩张}沿 $d_0$ 的\Term{inverse-image}{逆像}，并要求 $\varphi$ 满足一个与 $d_0d_1=0$ 相容的条件，下面将具体说明它。而 $d_0$ 的定义 (3.5.2) 立即给出
\[
E'=\pi^*(E)\operatorname{pr}_1^*(E)^{-1}\operatorname{pr}_2^*(E)^{-1},
\]
这个\Term{torsor}{挠子}的一个\Term{trivialization}{平凡化}可以解释为一个\Term{isomorphism}{同构}（也记为 $\varphi$）
\[
\pi^*(E)\xleftarrow[\sim]{\varphi}\operatorname{pr}_1^*(E)\operatorname{pr}_2^*(E).
\]
另一方面，具体写出 $\varphi$ 与 $d_1d_0=0$ 的相容条件，便发现它由 (1.2.1) 的\Term{commutativity}{交换性}条件和 (1.1.4.1) 的\Term{associativity}{结合性}条件表达；这两个条件分别对应于 $L_2$ 的两个因子。此时，(*) 是\Term{equivalence-of-categories}{范畴等价}这一事实，正是 1.1.6 中建立的对应。这就完成了 3.5.3 的证明。
\end{proof}

\textbf{注记 3.5.4。} a) 若能构造 $P$ 的一个关于 $P$ 具有\Term{functoriality}{函子性}的\Term{resolution}{消解} $L_\bullet(P)$，使其各\Term{set-component}{分量}（关于 $P$ 具有\Term{functoriality}{函子性}地）\Term{isomorphic}{同构}于由 $P$ 的\Term{cartesian-power}{笛卡儿幂}生成的\Term{free-module}{自由} $\mathbb Z$-\Term{module}{模}的\Term{direct-sum}{直和}，从而加强 3.5.3，将是很有意义的。
\SourcePage{211}
事实上，考虑到 (1.4.5) 型的\Term{isomorphism}{同构}，这将给出一种部分计算 $\operatorname{Ext}^i(P,A)$ 的方法：把它视为一个\Term{spectral-sequence}{谱序列}的\Term{abutment}{收敛项}，该\Term{spectral-sequence}{谱序列}的初始项可用 $H^j(P\times P\cdots\times P,A)$ 具体表达。（这里，我们只得到关于 $\operatorname{Ext}^1$ 的很不完整的信息，而不涉及更高的 $\operatorname{Ext}^i$。）无论如何，有一种自然方法可将部分\Term{resolution}{消解} $L_\bullet(P)$ 再向前延长一步。\footnote{L. BREEN 已构造了这样的\Term{resolution}{消解}（校订者注）。}

b) 可以在 $A$-\Term{module}{模}的\Term{category}{范畴}中给出构造 (3.5.1) 和 3.5.3 的一个变式，其中 $A$ 是 $\mathcal T$ 的一个固定的\Term{ring}{环}。由于必须使 1.1.6 更复杂（参见 2.10.3），\Term{complex}{复形} $L_\bullet(P)$ 的\Term{form}{形式}也更复杂。得到
\[
\left\{\begin{aligned}
L_0&=A[P],\\
L_1&=A[P\times P]+A[A\times P],\\
L_2&=A[P\times P\times P]+A[P\times P]+A[A\times A\times P]+A[A\times P\times P],
\end{aligned}\right.\tag{3.5.5}
\]
其向 $P$ \Term{augmentation}{增广}的\Term{complex}{复形}结构由公式 (3.5.2) 加上下面的公式给出：
\[
\left\{\begin{aligned}
d_0[\lambda,p]&=[\lambda p]-\lambda[p],\\
d_1[\lambda,\mu,p]&=[\lambda,\mu p]-[\lambda\mu,p]+\lambda[\mu,p],\\
d_1[\lambda,p,q]&=[\lambda,p+q]-[\lambda,p]-[\lambda,q]+\lambda[p,q]-[\lambda p,\lambda q].
\end{aligned}\right.\tag{3.5.6}
\]
关系 $d_0d_1=0$ 来自以下关系：一是将\Term{differential}{算子} $d_0,d_1$ \Term{restriction}{限制}到 (3.5.2) 中所考虑的\Term{set-component}{分量}时的类似关系（如前所见，它\Term{equivalence}{等价}于 $P$ 上给定加法律的\Term{associativity}{结合性}与\Term{commutativity}{交换性}）；二是关系 $d_0d_1(\lambda,\mu,p)=0$ 和 $d_0d_1(\lambda,p,q)=0$。立即可验证，后两个关系分别\Term{equivalence}{等价}于
\SourcePage{212}
$A$ 在 $P$ 上给定\Term{action}{作用}的\Term{associativity}{结合性}（即 $\lambda(\mu p)=(\lambda\mu)p$）和它关于加法的\Term{distributivity}{分配性}（即 $\lambda(p+q)=\lambda p+\lambda q$）。对一个 $G_P$-\Term{torsor}{挠子} $E$，给定 $d_0^*(E)$ 的一个\Term{trivialization}{平凡化}，\Term{equivalence}{等价}于给定 $E$ 上的一个内\Term{composition-law}{合成律}（满足前面已经具体说明的、与 $E$ 的内\Term{composition-law}{合成律}及\Term{torsor}{挠子}结构相容的条件），以及 $A$ 在 $E$ 上的一个外\Term{composition-law}{合成律}，后者与 $A$ 在 $P$ 和 $G$ 上的\Term{action}{作用}相容。此\Term{trivialization}{平凡化}与 $d_1^*d_0^*(E)$ 的显然\Term{trivialization}{平凡化}相容，便\Term{equivalence}{等价}于前面已经具体说明的条件（表达 $E$ 的内\Term{composition-law}{合成律}是一个\Term{commutative-group}{交换群}的律，且给出 $P$ 以 $G$ 为核的\Term{extension}{扩张}），再加上外律的\Term{associativity}{结合性}及其关于内律的\Term{distributivity}{分配性}；这些条件\Term{equivalence}{等价}于说，$E$（对于这两个\Term{composition-law}{合成律}及其\Term{torsor}{挠子}结构而言）是 $A$-\Term{module}{模} $P$ 以 $G$ 为核的一个 $A$-\Term{module-extension}{模扩张}。由此，基本上用与 3.5.3 相同的方法，可得 $L_\bullet$ 是\Term{module}{模} $P$ 的一个部分\Term{resolution}{消解}（即 $H_1(L_\bullet)=0$）。当然，若能把这个部分\Term{resolution}{消解}延长为类似类型的无限\Term{resolution}{消解}，使各\Term{set-component}{分量}都是形如 $A(A\times A\cdots\times P\times P\cdots)$ 的\Term{module}{模}，也仍然是很有意义的；比较 a) 中\Term{commutative-group}{交换群}情形的类似问题。

\SourcePage{213}
\SGAsection{3.6}{应用：描述给定交换群对 $(P,Q)$ 的双扩张的余纤维范畴}
设 $P,Q$ 为\Term{topos}{拓扑斯} $\mathcal T$ 的两个 $\mathbb Z$-\Term{module}{模}。任意选择 $P$ 的一个\Term{flat-resolution}{平坦消解} $L'_\bullet(P)$，使它在\Term{degree}{次数} $\leq2$ 处与 3.5 的\Term{augmented-complex}{增广复形} $L_\bullet(P)$ 相同——由 3.5.3，以及 $\mathcal T$ 上每个 $\mathbb Z$-\Term{module}{模}都是一个\Term{free-module}{自由}、因而\Term{flat}{平坦}的 $\mathbb Z$-\Term{module}{模}的商这一事实，这是可能的。类似地，选择 $Q$ 的一个\Term{flat-resolution}{平坦消解} $L'_\bullet(Q)$，使它在\Term{degree}{次数} $\leq2$ 处与 $L_\bullet(Q)$ 相同。因而有\Term{canonical}{典范}\Term{homomorphism}{同态}
\[
L_\bullet(P)\longrightarrow L'_\bullet(P),\qquad
L_\bullet(Q)\longrightarrow L'_\bullet(Q),\tag{3.6.1}
\]
诱导一个\Term{canonical}{典范}\Term{homomorphism}{同态}
\[
\begin{aligned}
L_\bullet(P,Q)&\overset{\mathrm{def}}=L_\bullet(P)\otimes L_\bullet(Q)\\
&\longrightarrow L'_\bullet(P,Q)\overset{\mathrm{def}}=L'_\bullet(P)\otimes L'_\bullet(Q)
\simeq P\mathbin{\otimes^{\mathrm L}}Q,
\end{aligned}\tag{3.6.2}
\]
它在\Term{degree}{次数} $\leq2$ 的\Term{set-component}{分量}上是\Term{isomorphism}{同构}。因此由定义有
\[
\Psi_{L_\bullet(P,Q)}=\Psi_{L_\bullet(P',Q')}.\tag{3.6.3}
\]

\begin{SGAresult}{\Term{theorem}{定理}}{3.6.4}
在\Term{topos}{拓扑斯} $\mathcal T$ 的\Term{commutative-group}{交换群}\Term{category}{范畴} $\mathcal C=\operatorname{GRAB}(\mathcal T)$ 上，由 $(P,Q)$ 以变动的\Term{abelian-group}{阿贝尔群} $G$ 为核的\Term{biextension}{双扩张}组成的\Term{additive}{加性}\Term{cofibered-category}{余纤维范畴}（2.4），\Term{canonically-equivalent}{典范等价}于按 3.1 由 $L_\bullet(P,Q)=L_\bullet(P)\otimes L_\bullet(Q)$ 定义的\Term{additive}{加性}\Term{cofibered-category}{余纤维范畴} $\Psi_{L_\bullet(P,Q)}=\Psi_{L'_\bullet(P,Q)}$。
\end{SGAresult}
由此及 3.2.5，得到本节前面已经宣布的主要结果：
\SourcePage{214}
\begin{SGAresult}{推论}{3.6.5}
固定 $P,Q$ 后，有关于 $G$ 具有\Term{functoriality}{函子性}的\Term{canonical-isomorphism}{典范同构}
\[
\operatorname{Biext}^1(P,Q;G)\simeq\operatorname{Ext}^1(P\mathbin{\otimes^{\mathrm L}}Q,G).\tag{3.6.5.1}
\]
\end{SGAresult}
\begin{proof}
下面证明 3.6.5；为此，先具体写出 $L_\bullet(P,Q)$ 在\Term{degree}{次数} $0,1,2$ 的\Term{set-component}{分量}：
\[
\left\{\begin{aligned}
L_0(P,Q)&=L_0(P)\otimes L_0(Q)=\mathbb Z[P\times Q],\\
L_1(P,Q)&=L_0(P)\otimes L_1(Q)+L_1(P)\otimes L_0(Q)\\
&=\mathbb Z[P\times Q\times Q]+\mathbb Z[P\times P\times Q],\\
L_2(P,Q)&=L_0(P)\otimes L_2(Q)+L_1(P)\otimes L_1(Q)+L_2(P)\otimes L_0(Q)\\
&=\mathbb Z[P\times Q\times Q]+\mathbb Z[P\times Q\times Q\times Q]\\
&\quad+\mathbb Z[P\times P\times Q\times Q]\\
&\quad+\mathbb Z[P\times P\times Q]+\mathbb Z[P\times P\times P\times Q].
\end{aligned}\right.\tag{3.6.6}
\]
为通过\Term{fiber-category}{纤维范畴} $\Psi_{L_\bullet(P,Q)}(G)$ 具体描述\Term{cofibered-category}{余纤维范畴} $\Psi_{L_\bullet(P,Q)}$，利用 1.4，把 $L_i(P,Q)$（$i=0,1,2$）以 $G$ 为核的\Term{extension}{扩张}用\Term{torsor}{挠子}具体表达。这样，$L_0(P\times Q)$ 以 $G$ 为核的一个\Term{extension}{扩张}与一个 $G_{P\times Q}$-\Term{torsor}{挠子}等同；$L_1(P,Q)$ 以 $G$ 为核的一个\Term{extension}{扩张}与一个 $G_{P\times Q\times Q}$-\Term{torsor}{挠子}和一个 $G_{P\times P\times Q}$-\Term{torsor}{挠子}构成的二元组等同；类似地，$L_2(P,Q)$ 以 $G$ 为核的一个\Term{extension}{扩张}与五个\Term{torsor}{挠子}构成的系统等同，其结构\Term{group}{群}分别由 $G$ 向底 $P\times Q\times Q$、$P\times Q\times Q\times Q$、$P\times P\times Q\times Q$、$P\times P\times Q$、$P\times P\times P\times Q$ 作\Term{base-change}{基变换}得到。利用这个对应，可将由 $L_0(P,Q)$ 以 $G$ 为核的一个\Term{extension}{扩张} $X$，及 $L_1(P,Q)$ 上 $d_0^*(X)$ 的一个\Term{trivialization}{平凡化} $\alpha$ 所成的二元组 $(X,\alpha)$，解释为 $P\times Q$ 上的一个\Term{torsor}{挠子} $E$，连同一对\Term{trivialization}{平凡化} $(\varphi,\psi)$；它们分别\Term{trivialization}{平凡化} $P\times Q\times Q$ 和 $P\times P\times Q$ 上的\Term{torsor}{挠子}对 $(E_1,E_2)$，后者由 $E$ 借助
\SourcePage{215}
$d_0:L_1(P,Q)\to L_0(P,Q)$ 取\Term{inverse-image}{逆像}得到。具体写出\Term{differential}{微分算子} $d_0$，如同在 3.5.3 的证明中一样，可知 $\varphi,\psi$ \Term{solvable}{可解}释为\Term{isomorphism}{同构}
\[
\varphi:\pi_Q^*(E)\xleftarrow{\sim}q_1^*(E)q_2^*(E),\qquad
\psi:\pi_P^*(E)\xleftarrow{\sim}p_1^*(E)p_2^*(E),
\]
其中 $\pi_Q,q_1,q_2$ 是从 $P\times Q\times Q$ 到 $Q$ 的三个\Term{morphism}{态射}，分别由 $Q\times Q\to Q$ 的三个\Term{morphism}{态射}——\Term{composition-law}{合成律}、第一\Term{projection}{投影}和第二\Term{projection}{投影}——得到，而 $\pi_P,p_1,p_2:P\times P\times Q\to P\times Q$ 以\Term{symmetric}{对称}方式定义。因此，$\varphi,\psi$ \Term{solvable}{可解}释为 2.1 中出现的 (2.0.4) 和 (2.0.7) 型数据。最后，二元组 $(X,\alpha)$ 确为 $\Psi_{L_\bullet(P,Q)}(G)$ 的\Term{object}{对象}这一条件，即部分\Term{trivialization}{平凡化} $\alpha$ 与关系 $d_0d_1=0$ 相容，针对结构 $(E,\varphi,\psi)$ 而言，表达为五个相容关系；它们分别用 $P\times Q\times Q$、$P\times Q\times Q\times Q$、$P\times P\times Q\times Q$、$P\times P\times Q$、$P\times P\times P\times Q$ 上的\Term{torsor}{挠子}表达，对应于 $L_2(P,Q)$ 的表达式 (3.6.6) 中的五项。如同 3.5.3 的证明中一样可验证，前两个关系正是\Term{commutativity}{交换性}关系 (2.0.6) 和\Term{associativity}{结合性}关系 (2.0.5)，两者合在一起表达 $\varphi$ 在 $E$ 上定义了 $Q_P$ 以 $G_P$ 为核的一个\Term{commutative-group-extension}{交换扩张}结构。\Term{symmetric}{对称}地，最后两个条件表达 $\psi$ 在 $E$ 上定义了 $P_Q$ 以 $G_P$ 为核的一个\Term{commutative-group-extension}{交换扩张}结构。最后，中间的条件涉及 $P\times P\times Q\times Q$ 上的\Term{torsor}{挠子}，它正是图 (2.1.1) 的\Term{commutativity}{交换性}；后者表达 $\varphi$ 与 $\psi$ 相容，即（考虑到其余条件）它们在 $E$ 上定义了 $(P,Q)$ 以 $G$ 为核的\Term{biextension}{双扩张}结构。这些解释的验证
\SourcePage{216}
基本上是显然的，留给读者。读者想必会与我们一样承认，前面的论述确实定义了 $\mathcal C=\operatorname{GRAB}(\mathcal T)$ 上的一个\Term{cofibered-category}{余纤维范畴}\Term{equivalence-of-categories}{等价}，它就是 3.6.4 中所宣布的\Term{equivalence-of-categories}{等价}。
\end{proof}

\textbf{注记 3.6.7。} 2.10 中指出的\Term{biextension}{双扩张}概念的三种推广，每一种都给出 3.6.4（因而也给出 3.6.5）的相应变式。推广到任意多个因子的情形（参见 2.10.2）是显然的；而一旦在 $A$-\Term{module}{模}的情形得到了 3.5.4 b) 中所述的 $L_\bullet(P)$ 构造的变式，推广到 $A$-\Term{module}{模}的\Term{biextension}{双扩张}情形（参见 2.10.3）也是直接的。从\Term{homological-algebra}{同调代数}观点看，推广到不假定 $P,Q$ 交换的情形（参见 2.10.1）似乎更有意思。这时应给 $\mathcal T$ 的一个（未必交换的）\Term{group}{群} $P$ 对应一个二阶\Term{truncation}{截断}的\Term{complex}{复形} $K_\bullet(P)$；它的描述与 3.5 中 $L_\bullet(P)$ 的描述类似，区别在于取 $K_2(P)=\mathbb Z[P\times P\times P]$（即略去 $L_\bullet(P)$ 中出现的第一个因子），而 $d_0,d_1$ 仍按 (3.5.2) 定义。这里 $H_0(K_\bullet(P))$ 与由 $P$ \Term{abelianization}{阿贝尔化}得到的\Term{group}{群} $P/[P,P]$ 等同，也可与 $P$ 以 $\mathbb Z$ 为系数的第 $1$ 次\Term{homology-group}{同调群}所成的\Term{sheaf}{层} $\underline H_1(P,\mathbb Z)$ 等同；同样，$H_1(K_\bullet(P))$ \Term{isomorphic}{同构}于\Term{sheaf}{层} $\underline H_2(P,\mathbb Z)$。事实上，$K_\bullet(P)$ 正是 $P$ 以 $\mathbb Z$ 为系数的\Term{chain-complex}{链复形} $C_\bullet(P,\mathbb Z)$ 的二阶\Term{truncation}{截断}。于是，建立 3.5.3 的论证也同样证明，$\mathcal C=\operatorname{GRAB}(\mathcal T)$ 上的\Term{cofibered-category}{余纤维范畴} $\Psi_{K_\bullet(P)}$
\SourcePage{217}
\Term{canonically-equivalent}{典范等价}于 $P$ 以变动的\Term{commutative-group}{交换群} $G$ 为核的\Term{central-extension}{中心扩张}所成的\Term{category}{范畴}。同样，对于两个\Term{group}{群} $P,Q$，可得到 2.10.1 意义下 $(P,Q)$ 的\Term{biextension}{双扩张}的\Term{cofibered-category}{余纤维范畴}与\Term{cofibered-category}{余纤维范畴} $\Psi_{K_\bullet(P,Q)}$ 的一个\Term{equivalence-of-categories}{等价}，其中
\[
K_\bullet(P,Q)=K_\bullet(P)\mathbin{\otimes^{\mathrm L}}K_\bullet(Q).
\]
这些思考提示：当 $P$ 交换时，也许可以对\Term{chain-complex}{链复形} $C_\bullet(P,\mathbb Z)$ 作适当修改，从而找到 3.5.4 a) 中所希望的 $P$ 的一个\Term{canonical}{典范}\Term{resolution}{消解}。

\SGAsection{3.7}{各种相容性}
本小节考察 $P,Q$ 变动时\Term{cofibered-category}{余纤维范畴}的\Term{equivalence-of-categories}{等价} 3.6.4 的行为。$(P,Q)$ 以变动的 $G$ 为核的\Term{biextension}{双扩张}的\Term{cofibered-category}{余纤维范畴}，关于变动的 $P,Q$ 的\Term{variance}{变性}，原则上已在 2.4 中定义，也可表述为存在一个\Term{cofibered-functor}{余纤维化函子} (2.6.1)。注意，关于 $P,Q$ 的\Term{variance}{变性}是\Term{contravariant}{反变的}。\Term{cofibered-category}{余纤维范畴} $\Psi_{L_\bullet(P,Q)}$ 关于变动的 $P,Q$ 的\Term{variance}{变性}也很清楚，因为 $L_\bullet(P,Q)=L_\bullet(P)\otimes L_\bullet(Q)$ 关于二元组 $(P,Q)$ 具有\Term{functoriality}{函子性}，而 $\Psi_{L_\bullet}$ 关于 $L_\bullet$ 是\Term{contravariant}{反变的}（3.4）。于是第一个相容性可以表述为：对每个\Term{morphism}{态射}三元组
\[
P'\longrightarrow P,\qquad Q'\longrightarrow Q,\qquad G\longrightarrow G',\tag{3.7.1}
\]
下列\Term{functor}{函子}图
\[
\begin{tikzcd}[column sep=large]
\Psi_{L_\bullet(P,Q)}(G)\arrow[r]\arrow[d]&\Psi_{L_\bullet(P',Q')}(G')\arrow[d]\\
\operatorname{BIEXT}(P,Q;G)\arrow[r]&\operatorname{BIEXT}(P',Q';G')
\end{tikzcd}\tag{3.7.2}
\]
\SourcePage{218}
在一个\Term{canonical-isomorphism}{典范同构}意义下交换。我们承认这种相容性，因为具体写出交换\Term{isomorphism}{同构}虽很显然，却同样繁琐（而且它所要求的\Term{cofibered-category}{余纤维范畴}\Term{equivalence-of-categories}{等价} 3.6.4 的完整具体描述也是如此）。

\textbf{3.7.3。} 上述相容性特别说明，\Term{isomorphism}{同构} (3.6.5.1) 不仅关于 $G$，而且关于 $P,Q$ 都具有\Term{functoriality}{函子性}。同样，而且更直接地，对借助 (3.1.8) 定义的下列\Term{isomorphism}{同构}，也可以陈述类似的\Term{functoriality}{函子性}，正如 (3.6.5.1) 是借助 (3.2.4) 定义的一样：
\[
\operatorname{Biext}^0(P,Q;G)\xrightarrow{\sim}
\operatorname{Ext}^0(P\mathbin{\otimes^{\mathrm L}}Q,G)\simeq\operatorname{Hom}(P\otimes Q,G).\tag{3.7.4}
\]
但应注意，我们在 (2.5.4.1) 中已经得到过这样的\Term{isomorphism}{同构}，并在 2.6 中指出它关于三个变元都具有\Term{functoriality}{函子性}。而立即可验证，两个\Term{isomorphism}{同构} (3.7.4) 和 (2.5.4.1) 是相同的。

\textbf{3.7.5。} \Term{isomorphism}{同构} (3.6.5) 的意义在于：它把具有几何性质的一个重要\Term{group}{群} $\operatorname{Biext}^1(P,Q;G)$，用一个很简单的\Term{homological-algebra}{同调代数}\Term{invariant}{不变量} $\operatorname{Ext}^1(P\mathbin{\otimes^{\mathrm L}}Q,G)$ 来表达；后者置于\Term{invariant}{不变量} $\operatorname{Ext}^i(P\mathbin{\otimes^{\mathrm L}}Q,G)$ 的序列之中，可以对这一序列应用通常的\Term{homological-algebra}{同调代数}\Term{form}{形式}体系。特别地，这些 $\operatorname{Ext}^i$ 组成的序列，关于 $P,Q,G$ 三个变元中的每一个，都是 [3] 意义下的\Term{cohomological-functor}{上同调函子}（或按 [5] 的术语，称为\Term{connected-sequence-of-functors}{函子的相联系列}）。特别地，一个\Term{exact-sequence}{正合列}
\[
0\longrightarrow G'\longrightarrow G\longrightarrow G''\longrightarrow0\tag{3.7.5.1}
\]
\SourcePage{219}
（或相应地，
\[
0\longrightarrow P'\longrightarrow P\longrightarrow P''\longrightarrow0,
\quad\text{或 }0\longrightarrow Q'\longrightarrow Q\longrightarrow Q''\longrightarrow0\tag{3.7.5.2}
\]
）给出 $\operatorname{Ext}^i$ 的一个无限\Term{exact-sequence}{正合列}，其前六项（只涉及 $\operatorname{Ext}^0$ 和 $\operatorname{Ext}^1$）借助 3.6.5 和 3.7.4 可用\Term{biextension}{双扩张}解释。我们已经在\Term{left-exact}{左正合}\Term{cofibered-category}{余纤维范畴}的概念框架内（3.3.6），指出了由 (3.7.5.1) \Term{derive}{导出}的\Term{exact-sequence}{正合列}，尤其是\Term{coboundary-operator}{上边界算子}
\[
\partial:\operatorname{Biext}^0(P,Q;G'')\longrightarrow\operatorname{Biext}^1(P,Q;G')\tag{3.7.5.3}
\]
的几何解释。可以尝试对\Term{short-exact-sequence}{短正合列} (3.7.5.2) 之一，比如第一个，所对应的六项\Term{exact-sequence}{正合列}给出类似解释；考虑到 3.7.3，这\Term{equivalence}{等价}于给\Term{coboundary-homomorphism}{上边界同态}
\[
\partial:\operatorname{Biext}^0(P',Q;G)\longrightarrow\operatorname{Biext}^1(P'',Q;G)\tag{3.7.5.4}
\]
作几何解释。事实上有：

\begin{SGAresult}{命题}{3.7.6}
在 $\operatorname{GRAB}(\mathcal T)^\circ\times\operatorname{GRAB}(\mathcal T)^\circ\times\operatorname{GRAB}(\mathcal T)$ 上，由变动的 $P,Q,G$ 所给出的 $(P,Q)$ 以 $G$ 为核的\Term{biextension}{双扩张}所成的\Term{cofibered-category}{余纤维范畴}（参见 (2.6.1)），不仅关于变元 $G$ \Term{left-exact}{左正合}，而且关于变元 $P,Q$ 中的每一个也\Term{left-exact}{左正合}。此外，由 (3.7.5.2) 对应的 $\operatorname{Ext}^i$ 的\Term{exact-sequence}{正合列}（借助\Term{isomorphism}{同构} 3.6.5 和 3.7.4）所得的\Term{coboundary-homomorphism}{上边界同态} (3.7.5.4)，就是 3.3.6 中回顾的、$\mathcal C=\operatorname{GRAB}(\mathcal T)^\circ$ 上\Term{left-exact}{左正合}\Term{cofibered-category}{余纤维范畴}理论的\Term{coboundary-homomorphism}{上边界同态}，
\SourcePage{220}
应用于 $Q,G$ 固定、$P$ 变动时 $(P,Q)$ 以 $G$ 为核的\Term{biextension}{双扩张}的\Term{cofibered-category}{余纤维范畴}（须保留 3.4.9 中可能有\Term{signature}{符号}错误这一点，参见 3.4.9.2）。
\end{SGAresult}
\begin{proof}
下面勾勒 3.7.6 的证明。借助 3.6.4 及相容性 (3.7.2)，将 $(P,Q)$ 以 $G$ 为核的\Term{biextension}{双扩张}解释为 $\Psi_{L_\bullet(P,Q)}(G)$ 的\Term{object}{对象}。考虑由\Term{augmentation}{增广} $L_\bullet(P)\to P$ 定义的\Term{homomorphism}{同态}
\[
L_\bullet(P,Q)=L_\bullet(P)\otimes L_\bullet(Q)\longrightarrow P\otimes L_\bullet(Q).
\]
由 3.5.3 和 $L_\bullet(Q)$ 的\Term{flatness}{平坦性}，此\Term{homomorphism}{同态}在\Term{object}{对象} $H_0,H_1$ 上诱导\Term{isomorphism}{同构}，因而由 3.4.7 诱导一个\Term{equivalence-of-categories}{等价}
\[
\Psi_{P\otimes L_\bullet(Q)}\longrightarrow\Psi_{L_\bullet(P,Q)}.
\]
这使我们能将 $(P,Q)$ 以 $G$ 为核的\Term{biextension}{双扩张}解释为 $\Psi_{P\otimes L_\bullet(Q)}(G)$ 的\Term{object}{对象}。此时，\Term{exact-sequence}{正合列} (3.7.5.2) 给出一个\Term{chain-complex}{链复形}的\Term{exact-sequence}{正合列}
\[
0\longrightarrow P'\otimes L_\bullet(Q)\longrightarrow P\otimes L_\bullet(Q)
\longrightarrow P''\otimes L_\bullet(Q)\longrightarrow0.\tag{3.7.6.1}
\]
将 3.4.8 应用于这个\Term{complex}{复形}的\Term{exact-sequence}{正合列}，就得到所宣布的\Term{left-exact}{左正合性}。（直接验证也同样是立即的。）另一方面，\Term{homomorphism}{同态} (3.7.5.4) 正是由这个\Term{exact-sequence}{正合列}取以 $G$ 为值的 $\operatorname{Ext}^i$ 后得到的\Term{coboundary-homomorphism}{上边界同态} $\operatorname{Ext}^0\to\operatorname{Ext}^1$；因此，3.7.6 的最后一个断言是 3.4.9 的特例。
\end{proof}

\SourcePage{221}
\SGAsection{3.8}{在 $P$ 和 $Q$ 的扩张上平凡化的 $(P,Q)$ 的双扩张}
先考虑两个\Term{abelian-group}{阿贝尔群}的\Term{exact-sequence}{正合列}
\[
0\longrightarrow L_1\xrightarrow{i}L_0\xrightarrow{u}P\longrightarrow0,
\qquad0\longrightarrow M_1\xrightarrow{j}M_0\xrightarrow{v}Q\longrightarrow0.\tag{3.8.1}
\]
我们要描述 $(P,Q)$ 以\Term{abelian-group}{阿贝尔群} $G$ 为核的\Term{biextension}{双扩张} $E$ 所成的\Term{category}{范畴}，其中要求 $E$ 沿 $(u,v)$ 的\Term{inverse-image}{逆像} $(u,v)^*(E)$ 是\Term{trivial-biextension}{平凡的}；也就是说，存在一个\Term{homomorphism}{同态} $s$ 使下面的三角形交换，
\[
\begin{tikzcd}[column sep=large]
E\arrow[d]&\\
P\times Q&L_0\times M_0\arrow[l,"u\times v"']\arrow[ul,"s"']
\end{tikzcd}\tag{3.8.2}
\]
并且 $s$ 对两个变元中的每一个都满足所需的\Term{additivity}{加性}条件。因此，给定这样的 $s$，就可以将 $(u,v)^*(E)$ 与 $(L_0,M_0)$ 以 $G$ 为核的\Term{trivial-biextension}{平凡双扩张} $E_0=L_0\times M_0\times G$ 等同。对于固定的 $G$，$(P,Q)$ 以 $G$ 为核的\Term{biextension}{双扩张}的\Term{cofibered-category}{余纤维范畴}关于 $P,Q$ 的\Term{left-exact}{左正合性}（3.7.6）此时表明：给定 $(P,Q)$ 的\Term{biextension}{双扩张} $E$ 及其\Term{inverse-image}{逆像}的一个\Term{trivialization}{平凡化} $s$，基本上\Term{equivalence}{等价}于给定 $(L_0,M_0)$ 以 $G$ 为核的\Term{trivial-biextension}{平凡双扩张} $E_0$ 的两个部分\Term{trivialization}{平凡化}，它们分别定义在 $(L_1,M_0)$ 和 $(L_0,M_1)$ 上，并且在 $(L_1,M_1)$ 上相同。由于 $E_0$ 在 $(L_1,M_0)$、$(L_0,M_1)$ 和 $(L_1,M_1)$ 上诱导的\Term{biextension}{双扩张}可与\Term{trivial-biextension}{平凡双扩张}等同，因此给定 $(L_1,M_0)$ 上的部分\Term{trivialization}{平凡化}（相应地，$(L_0,M_1)$ 上的部分\Term{trivialization}{平凡化}），\Term{equivalence}{等价}于给定一个\Term{homomorphism}{同态}（或\Term{pairing}{配对}）
\[
f:L_1\otimes M_0\longrightarrow G
\qquad\text{（相应地，$g:L_0\otimes M_1\longrightarrow G$）},\tag{3.8.3}
\]
\SourcePage{222}
而所考虑的相容条件只表示这两个\Term{homomorphism}{同态}在 $L_1\otimes M_1$ 上相同，即对 $L_1,M_1$ 在 $\mathcal T$ 的某个\Term{object}{对象}上的\Term{section}{截面} $x_1,y_1$，有
\[
f(x_1\otimes y_1)=g(x_1\otimes y_1).\tag{3.8.4}
\]
用 (3.8.2) 的\Term{homomorphism}{同态} $s$ 表达，\Term{form}{型} $f,g$ 由以下条件刻画：
\[
s(x_1,y_0)=f(x_1\otimes y_0)e_{E/Q}(y_0),\qquad
s(x_0,y_1)=g(x_0\otimes y_1)e_{E/P}(x_0),\tag{3.8.5}
\]
这再次显示了条件 (3.8.4) 的必要性。

选定\Term{trivialization}{平凡化} $s$ 后，其余的\Term{trivialization}{平凡化}都是形如下式的\Term{morphism}{态射} $s':L_0\times M_0\to E$：
\[
s'(x_0,y_0)=s(x_0,y_0)+h(x_0\otimes y_0),\tag{3.8.6}
\]
其中 $h:L_0\otimes M_0\to G$ 是任意的\Term{form}{型}。于是从 (3.8.5) 清楚可见，由 $s$ 定义的\Term{form}{型} $f,g$ 被改成\Term{form}{型}
\[
\begin{aligned}
f'(x_1\otimes y_0)&=f(x_1\otimes y_0)+h(x_1\otimes y_0),\\
g'(x_0\otimes y_1)&=g(x_0\otimes y_1)+h(x_0\otimes y_1).
\end{aligned}\tag{3.8.7}
\]
由此得到：

\begin{SGAresult}{命题}{3.8.8}
$\operatorname{Biext}^1(P,Q;G)$ 中由沿 $(u,v)$ 的\Term{inverse-image}{逆像}为 $(L_0,M_0)$ 的\Term{trivial-biextension}{平凡双扩张}的\Term{biextension}{双扩张}的类组成的\Term{subgroup}{子群}（即 $\operatorname{Biext}^1(P,Q;G)\to\operatorname{Biext}^1(L_0,M_0;G)$ 的\Term{kernel}{核}），借助上面的描述，\Term{canonically-isomorphic}{典范同构}于下述\Term{quotient-group}{商群}：
\SourcePage{223}
由满足相容条件 (3.8.4) 的 (3.8.3) 中\Term{form}{型}的二元组 $(f,g)$ 所成的\Term{group}{群}，模去那些由某个\Term{form}{型} $h:L_0\otimes M_0\to G$ \Term{restriction}{限制}得到的\Term{form}{型}的二元组。
\end{SGAresult}

现在，我们要在考虑 3.6.5 的\Term{canonical-isomorphism}{典范同构}
\[
\operatorname{Biext}^1(P,Q;G)\simeq\operatorname{Ext}^1(P\mathbin{\otimes^{\mathrm L}}Q,G)
\]
之后，用\Term{homological-algebra}{同调代数}解释上述\Term{isomorphism}{同构}。为此，把数据 (3.8.1) 解释为分别\Term{resolution}{消解} $P,Q$ 的两个\Term{complex}{复形}：
\[
L_\bullet:0\longrightarrow L_1\longrightarrow L_0\longrightarrow0,
\qquad M_\bullet:0\longrightarrow M_1\longrightarrow M_0\longrightarrow0,\tag{3.8.9}
\]
并考虑\Term{tensor-product}{张量积}\Term{complex}{复形} $L_\bullet\otimes M_\bullet$：
\[
0\longrightarrow L_1\otimes M_1\xrightarrow{d_2}
L_1\otimes M_0+L_0\otimes M_1\xrightarrow{d_1}L_0\otimes M_0\longrightarrow0.
\]
此外，我们知道在\Term{derived-category}{导出范畴}中有一个\Term{canonical}{典范}\Term{homomorphism}{同态}
\[
P\mathbin{\otimes^{\mathrm L}}Q\longrightarrow L_\bullet\otimes M_\bullet,
\]
由此得到一个同样\Term{canonical}{典范}的\Term{homomorphism}{同态}
\[
H^1\bigl(\operatorname{Hom}^\bullet(L_\bullet\otimes M_\bullet,G)\bigr)
\longrightarrow\operatorname{Ext}^1(P\mathbin{\otimes^{\mathrm L}}Q,G).\tag{3.8.10}
\]
还要注意，\Term{complex}{复形} $\operatorname{Hom}^\bullet(L_\bullet\otimes M_\bullet,G)$ 的 $1$-\Term{cocycle}{余圈}\Term{group}{群}，可与 (3.8.3) 中满足相容条件 (3.8.4) 的二元组 $(f,g)$ 所成的\Term{group}{群}等同：给这样的二元组对应\Term{set-component}{分量}为 $(f,g)$ 的\Term{homomorphism}{同态} $F:L_1\otimes M_0+L_0\otimes M_1\to G$，即
\[
F(x_1\otimes y_0+x_0\otimes y_1)=f(x_1\otimes y_0)+g(x_0\otimes y_1).\tag{3.8.11}
\]
此外，$0$-\Term{cochain}{上链} $h:L_0\otimes M_0\to G$ 的\Term{coboundary}{上边界}就是它在 $L_1\otimes M_0$ 和 $L_0\otimes M_1$ 上的\Term{restriction}{限制}所成的二元组，故 3.8.8 中所描述的二元组 $(f,g)$ 模去二元组 $(h,h)$ 后所得的\Term{group}{群}，
\SourcePage{224}
\Term{canonically-isomorphic}{典范同构}于 (3.8.10) 的左端。于是有下面的相容性：

\begin{SGAresult}{命题}{3.8.12}
若 $\xi\in\operatorname{Ker}(\operatorname{Biext}^1(P,Q;G)\to\operatorname{Biext}^1(L_0,M_0;G))$ 按 3.8.8 所述由二元组 $(f,g)$ 描述，则 $\xi$ 经\Term{isomorphism}{同构} (3.6.5.1) 在 $\operatorname{Ext}^1(P\mathbin{\otimes^{\mathrm L}}Q;G)$ 中的像，是由\Term{cocycle}{余圈} (3.8.11) 定义的左端元素经 (3.8.10) 所得的像。
\end{SGAresult}
这种相容性的验证留给读者。

\textbf{3.8.13。} 我们特别把这些结果应用于取
\[
L_0=\mathbb Z[P],\qquad M_0=\mathbb Z[Q]
\]
的情形，其中 $L_1,M_1$ 分别是\Term{canonical}{典范}\Term{homomorphism}{同态} $\mathbb Z[P]\to P$、$\mathbb Z[Q]\to Q$ 的\Term{kernel}{核}。注意，有\Term{canonical-isomorphism}{典范同构}
\[
L_0\mathbin{\otimes^{\mathrm L}}M_0\simeq L_0\otimes M_0\simeq\mathbb Z[P\times Q],
\]
因为 $L_0,M_0$ 都是\Term{flat}{平坦的}；由此借助 (1.4.5) 得到\Term{canonical-isomorphism}{典范同构}
\[
\operatorname{Ext}^i(L_0\mathbin{\otimes^{\mathrm L}}M_0,G)\simeq H^i(P\times Q,G_{P\times Q}).
\]
由此及 3.6.5，或者直接由定义及 1.4 可知：给每个 $(L_0,M_0)$ 以 $G$ 为核的\Term{biextension}{双扩张}对应它在 $P\times Q$ 上诱导的 $G_{P\times Q}$-\Term{torsor}{挠子}，所得\Term{functor}{函子}是\Term{equivalence-of-categories}{范畴等价}。（这个结论无需使用 $P,Q$ 上的\Term{additive}{加性}结构，仍然成立。）因此，若 $E$ 是 $(P,Q)$ 以 $G$ 为核的一个\Term{biextension}{双扩张}，则给定它在 $(L_0,M_0)$ 上的\Term{inverse-image}{逆像}的一个\Term{trivialization}{平凡化}，与给定仅视为 $P\times Q$ 上方\Term{object}{对象}的 $E$ 的一个\Term{section}{截面}，是同一件事；也就是说，与给定视为 $G_{P\times Q}$-\Term{torsor}{挠子}的 $E$ 的一个\Term{trivialization}{平凡化}相同。
\SourcePage{225}
带有这一附加结构的\Term{biextension}{双扩张}，因此由满足相容条件 (3.8.4) 的二元组 $(f,g)$ (3.8.3) 描述，而模\Term{isomorphism}{同构}的描述就是 3.8.8 中所给出的描述。

还要注意，采用 3.5 的记号，可将 $L_1$ 看作 $L_2(P)\to L_1(P)$ 的\Term{cokernel}{余核}，对 $M_1$ 也类似。由此，可将数据 $(f,g)$ 用一个\Term{homomorphism}{同态}系统具体表达：
\[
\left\{\begin{aligned}
f'&:L_1(P)\otimes L_0(Q)=\mathbb Z(P\times P\times Q)\longrightarrow G,\\
g'&:L_0(P)\otimes L_1(Q)=\mathbb Z(P\times Q\times Q)\longrightarrow G,
\end{aligned}\right.
\]
即\Term{set}{集合}\Term{sheaf}{层}的\Term{homomorphism}{同态}
\[
P\times P\times Q\longrightarrow G,\qquad P\times Q\times Q\longrightarrow G,
\]
满足五个相容条件（其中表达 $f',g'$ \Term{decomposable}{可分解}为 $f,g$ 的条件，考虑到 $L_2(P)$ 和 $L_2(Q)$ 各为两项之和的\Term{form}{形式}，需要写出 $2\times2$ 个方程；再加上条件 (3.8.4)）。这时可见，$f',g'$ 正是定义\Term{biextension}{双扩张}的部分\Term{multiplication}{乘法}数据，而上述五个条件正是 2.1 中考虑的\Term{biextension}{双扩张}的五条公理（参见 3.6.4 的证明）。

\SourcePage{226}
\SGAsection{}{参考文献}
\begin{enumerate}
\renewcommand{\labelenumi}{[\arabic{enumi}]}
\item M. Demazure、J. Giraud 和 M. Raynaud，《\Term{abelian-scheme}{阿贝尔概形}》（Schémas abéliens），Orsay 理学院讨论班，1967/68（待刊）。
\item J. Giraud，《\Term{nonabelian-cohomology}{非阿贝尔上同调}》（Cohomologie non abélienne），巴黎学位论文，1967（待刊）。
\item A. Grothendieck，《\Term{homological-algebra}{同调代数}的若干问题》（Sur quelques points d'Algèbre homologique）；Toh. Math. Journal，第 9 卷（1956），第 119--221 页。
\item A. Grothendieck，《\Term{additive}{加性}\Term{cofibred-category}{余纤维范畴}与\Term{relative-cotangent-complex}{相对余切复形}》（Catégories cofibrées additives et Complexe cotangent relatif），Lectures Notes in Math.，第 79 号，1968（Springer）。
\item J. L. Verdier，《\Term{abelian-category}{阿贝尔范畴}的\Term{derived-category}{导出范畴}》（Catégories dérivées des catégories abéliennes）（待刊）。
\item D. Mumford，《\Term{formal-group}{形式群}的\Term{biextension}{双扩张}》（Biextensions of formal groups），Tata Institute of Fundamental Research，1968（待刊）。
\end{enumerate}
